% =====================================================================
%  Git und GitHub – Ein praxisorientiertes Handbuch
%  Kompilieren: in VS Code speichern (LaTeX Workshop) oder  latexmk -pdf git-handbuch.tex
% =====================================================================
\documentclass[11pt, a4paper, parskip=half-, headsepline]{scrreprt}
% =====================================================================
%  Gemeinsame Präambel für alle Zusammenfassungen und Spickzettel
%  Einbinden mit:  \input{../../_latex/praeambel}
% =====================================================================

% --- Sprache und Schrift ---------------------------------------------
\usepackage[ngerman]{babel}
\usepackage[T1]{fontenc}
\usepackage{lmodern}
\usepackage{microtype}
% In Schreibmaschinenschrift keine Ligaturen: "--" bleibt zwei Striche, "<<" bleibt "<<"
\DisableLigatures{encoding = *, family = tt*}
\usepackage{csquotes}

% --- Seitenlayout (Ränder im Hauptdokument mit \geometry{} anpassen) --
\usepackage[a4paper]{geometry}

% --- Mathematik ------------------------------------------------------
\usepackage{amsmath, amssymb, mathtools}

% --- Tabellen und Listen ---------------------------------------------
\usepackage{array, booktabs, tabularx}
\usepackage{enumitem}

% --- Farben, Boxen, Grafiken -----------------------------------------
\usepackage[table]{xcolor}
\usepackage[most]{tcolorbox}
\usepackage{graphicx}
\graphicspath{{bilder/}}
\usepackage{tikz}
\usetikzlibrary{arrows.meta, positioning, calc}
\usepackage{forest}            % Bäume zeichnen (BST, Heap, AVL)

% --- Programmcode ----------------------------------------------------
\usepackage{listings}

% --- Links (immer zuletzt laden) -------------------------------------
\usepackage[hidelinks]{hyperref}

% =====================================================================
%  Farben
% =====================================================================
\definecolor{hsluBlau}{HTML}{1F4E79}
\definecolor{defFarbe}{HTML}{1F4E79}
\definecolor{merkFarbe}{HTML}{C55A11}
\definecolor{bspFarbe}{HTML}{548235}
\definecolor{pruefFarbe}{HTML}{A50021}
\definecolor{codeHintergrund}{HTML}{F5F5F5}

% =====================================================================
%  Boxen
%    \begin{definition}{Titel} ... \end{definition}
%    \begin{merke} ... \end{merke}
%    \begin{beispiel}{Titel} ... \end{beispiel}
%    \begin{pruefung} ... \end{pruefung}      (Prüfungstipp)
% =====================================================================
\tcbset{
  basis/.style={enhanced, breakable, boxrule=0.6pt, arc=2pt,
                left=6pt, right=6pt, top=4pt, bottom=4pt,
                fonttitle=\bfseries, before skip=8pt, after skip=8pt}
}
\newtcolorbox{definition}[1]{basis, colframe=defFarbe, colback=defFarbe!4,
  title={Definition: #1}}
\newtcolorbox{merke}{basis, colframe=merkFarbe, colback=merkFarbe!5,
  title={Merke}}
\newtcolorbox{beispiel}[1]{basis, colframe=bspFarbe, colback=bspFarbe!4,
  title={Beispiel: #1}}
\newtcolorbox{pruefung}{basis, colframe=pruefFarbe, colback=pruefFarbe!4,
  title={Prüfungstipp}}

% =====================================================================
%  Wochenziele-Tabelle
%    \begin{wochenziele}
%      \ziel{Lernpunkt}{Lernziel}
%    \end{wochenziele}
%  Die Nummern werden automatisch gezählt.
% =====================================================================
\newcounter{zielnr}
\newcommand{\ziel}[2]{\stepcounter{zielnr}\thezielnr & #1 & #2 \\}
\newenvironment{wochenziele}{%
  \setcounter{zielnr}{0}%
  \par\noindent
  \renewcommand{\arraystretch}{1.35}%
  \rowcolors{2}{hsluBlau!5}{white}%
  \tabularx{\linewidth}{@{}r>{\raggedright\arraybackslash}p{0.3\linewidth}>{\raggedright\arraybackslash}X@{}}
  \toprule
  \textbf{Nr.} & \textbf{Lernpunkt} & \textbf{Lernziel} \\
  \midrule
}{%
  \bottomrule
  \endtabularx
  \par
}

% =====================================================================
%  Programmcode
%    \begin{lstlisting}[language=Python] ... \end{lstlisting}
% =====================================================================
\lstset{
  basicstyle=\ttfamily\small,
  backgroundcolor=\color{codeHintergrund},
  keywordstyle=\color{hsluBlau}\bfseries,
  commentstyle=\color{gray}\itshape,
  stringstyle=\color{bspFarbe},
  numbers=left, numberstyle=\tiny\color{gray},
  frame=single, rulecolor=\color{gray!40},
  breaklines=true, showstringspaces=false, tabsize=4,
  literate={ä}{{\"a}}1 {ö}{{\"o}}1 {ü}{{\"u}}1
           {Ä}{{\"A}}1 {Ö}{{\"O}}1 {Ü}{{\"U}}1 {ß}{{\ss}}1
}

% =====================================================================
%  Kurzbefehle
% =====================================================================
\newcommand{\Oh}[1]{\mathcal{O}\!\left(#1\right)}   % \Oh{n^2}
\newcommand{\N}{\mathbb{N}}
\newcommand{\Z}{\mathbb{Z}}
\newcommand{\R}{\mathbb{R}}
\renewcommand{\vec}[1]{\mathbf{#1}}                 % Vektoren fett

\geometry{margin=2.4cm, top=2.6cm, bottom=2.8cm}
\AtBeginDocument{\shorthandoff{"}}            % " ohne Sonderbedeutung, wichtig für Befehle
\usetikzlibrary{decorations.pathreplacing}     % geschweifte Klammern in Grafiken

% --- Kopfzeile: Kapitelname links, Seitenzahl rechts -----------------
\usepackage[automark]{scrlayer-scrpage}
\pagestyle{scrheadings}
\clearpairofpagestyles
\ihead{\small\leftmark}
\ohead{\small\pagemark}
\setkomafont{pageheadfoot}{\normalfont\color{gray}}

% Inhaltsverzeichnis: Platz für zweistellige Kapitelnummern ("10.", "10.1.")
\RedeclareSectionCommand[tocnumwidth=2.0em]{chapter}
\RedeclareSectionCommand[tocnumwidth=2.8em, tocindent=2.0em]{section}

% =====================================================================
%  Bausteine
% =====================================================================

% Befehl im Fliesstext:  \cmd{git diff --cached}
\newcommand{\cmd}[1]{\texttt{\detokenize{#1}}}

% Abschnitt pro Befehl:  \befehl{git status}{Den Zustand abfragen}
\newcommand{\befehl}[2]{\section{\texorpdfstring{\texttt{\detokenize{#1}}}{#1}: #2}}

% Fachbegriff beim ersten Auftreten:  \begriff{Commit}
\newcommand{\begriff}[1]{\textbf{#1}}

% Kästen
\newtcolorbox{lernziele}{basis, colframe=hsluBlau, colback=white,
  title={Lernziele dieses Kapitels}}
\newtcolorbox{zusammenfassung}{basis, colframe=hsluBlau, colback=hsluBlau!4,
  title={Zusammenfassung}}
\newtcolorbox{praxis}{basis, colframe=bspFarbe, colback=bspFarbe!4,
  title={Aus der Praxis}}
\newtcolorbox{hintergrund}{basis, colframe=gray!65!black, colback=gray!5,
  title={Hintergrund}}
\newtcolorbox{erfahrung}{basis, colframe=merkFarbe, colback=merkFarbe!5,
  title={Erfahrungswert}}
\newtcolorbox{fehler}{basis, colframe=pruefFarbe, colback=pruefFarbe!4,
  title={Typischer Fehler}}
\newtcolorbox{achtung}{basis, colframe=pruefFarbe, colback=pruefFarbe!4,
  title={Achtung}}

% Tabelle "Teil | Bedeutung" zum Zerlegen eines Befehls
\newenvironment{teile}{%
  \par\noindent
  \renewcommand{\arraystretch}{1.3}%
  \tabularx{\linewidth}{@{}>{\ttfamily}l>{\raggedright\arraybackslash}X@{}}
  \toprule
  \normalfont\bfseries Teil & \bfseries Bedeutung \\
  \midrule
}{%
  \bottomrule
  \endtabularx
  \par
}

% Terminal: Befehle und Ausgaben (upquote: gerade ' und ` wie im Terminal)
\usepackage{upquote}
\lstset{upquote=true}
\lstdefinestyle{befehle}{language=bash, numbers=none, frame=single,
  backgroundcolor=\color{codeHintergrund},
  columns=fullflexible, keepspaces=true,
  alsoletter={-}, morekeywords={git, gh},
  deletekeywords={help, login, set, test, local, source, export},
  commentstyle=\color{gray}}
\lstdefinestyle{ausgabe}{language={}, numbers=none, frame=leftline,
  framerule=2pt, rulecolor=\color{gray!50}, backgroundcolor=\color{white},
  columns=fullflexible, keepspaces=true,
  basicstyle=\ttfamily\small\color{black!80}}

% Fehlermeldung mit Erklärung:  \meldung{Text}{Ursache}{Lösung}
% Im Text: \textquotesingle{} für ' und \{ \} \# \% für Sonderzeichen
\newcommand{\meldung}[3]{%
  \par\medskip
  {\raggedright\noindent{\color{pruefFarbe}\ttfamily\small #1}\par}
  \noindent\textbf{Ursache:} #2\par
  \noindent\textbf{Lösung:} #3\par
}

% =====================================================================
\subject{Handbuch}
\title{Git und GitHub}
\subtitle{Versionsverwaltung verstehen und anwenden,\\ wie sie in der Software-Entwicklung eingesetzt wird}
\author{Bahram Alizadah}
\date{\today}

\begin{document}
\maketitle
\tableofcontents

% ---------------------------------------------------------------------
\addchap{Vorwort}

Fast jede Software, die heute entsteht, wird mit \emph{Git} verwaltet: das Betriebssystem
auf dem Smartphone ebenso wie die App einer Bank oder das Analyseskript einer
Forschungsgruppe. Wer in der Informatik, in Data Science oder im Machine Learning arbeitet,
kommt daran nicht vorbei. In Stelleninseraten wird Git meist gar nicht mehr erwähnt, weil
es schlicht vorausgesetzt wird.

Dieses Handbuch richtet sich an Menschen, die Git \textbf{verstehen} wollen, nicht nur
bedienen. Es setzt keine Vorkenntnisse in Software-Entwicklung voraus. Jeder Fachbegriff
wird beim ersten Auftreten erklärt, und am Ende steht ein Glossar mit den deutschen und
englischen Bezeichnungen.

\section*{Aufbau}

\begin{description}[style=nextline, leftmargin=1.2cm]
  \item[Teil I: Grundlagen] erklärt, welches Problem Versionsverwaltung löst und wie Git
    intern denkt. Dieses \emph{mentale Modell} ist der Schlüssel: Wer es verstanden hat,
    kann sich fast jeden Befehl selbst herleiten.
  \item[Teil II: Arbeiten mit Git] führt durch Einrichtung, täglichen Arbeitsablauf und die
    wichtigsten Befehle, jeweils mit Zweck, Aufbau und Beispiel.
  \item[Teil III: Praxis] zeigt, wie in der Software-Industrie tatsächlich gearbeitet wird:
    gute Commits, Zusammenarbeit, Konflikte, Fehlerbehebung und Sicherheit.
  \item[Anhang] enthält eine Kurzreferenz zum Nachschlagen und das Glossar.
\end{description}

\section*{Konventionen}

\begin{itemize}
  \item Befehle und Dateinamen stehen in \cmd{Schreibmaschinenschrift}.
  \item Befehle, die im Terminal eingegeben werden, stehen in grauen Kästen.
        Alles nach \texttt{\#} ist ein Kommentar und wird nicht mit eingegeben:
\begin{lstlisting}[style=befehle]
git status                  # zeigt den aktuellen Zustand
\end{lstlisting}
  \item Ausgaben des Terminals sind mit einem grauen Balken am Rand markiert:
\begin{lstlisting}[style=ausgabe]
On branch main
nothing to commit, working tree clean
\end{lstlisting}
  \item Platzhalter, die durch eigene Werte ersetzt werden, stehen in spitzen Klammern:
        \cmd{git switch -c <name>} wird zu \cmd{git switch -c ads}.
\end{itemize}

Die farbigen Kästen haben feste Bedeutungen:

\begin{center}
\renewcommand{\arraystretch}{1.3}
\begin{tabularx}{\linewidth}{@{}l>{\raggedright\arraybackslash}X@{}}
  \toprule
  \textbf{Kasten} & \textbf{Inhalt} \\
  \midrule
  {\color{hsluBlau}\textbf{Lernziele / Zusammenfassung}} & was das Kapitel vermittelt \\
  {\color{defFarbe}\textbf{Definition}}                 & präzise Erklärung eines Fachbegriffs \\
  {\color{bspFarbe}\textbf{Aus der Praxis}}             & wie Teams in Unternehmen arbeiten und warum \\
  {\color{gray!65!black}\textbf{Hintergrund}}           & technische Vertiefung, beim ersten Lesen überspringbar \\
  {\color{merkFarbe}\textbf{Erfahrungswert}}            & was man meist erst nach Jahren lernt \\
  {\color{pruefFarbe}\textbf{Typischer Fehler / Achtung}} & häufige Stolpersteine und Gefahren \\
  \bottomrule
\end{tabularx}
\end{center}


\part{Grundlagen}
% ---------------------------------------------------------------------
\chapter{Warum Versionsverwaltung?}
\label{kap:warum}

\begin{lernziele}
\begin{itemize}[nosep]
  \item die Probleme kennen, die ohne Versionsverwaltung entstehen,
  \item verstehen, was ein Versionsverwaltungssystem leistet,
  \item Git und GitHub klar voneinander unterscheiden,
  \item wissen, warum Git in der Industrie zum Standard geworden ist.
\end{itemize}
\end{lernziele}

\section{Das Problem}

Jeder kennt Ordner, die so aussehen:

\begin{lstlisting}[style=ausgabe]
bericht.docx
bericht_v2.docx
bericht_v2_korrigiert.docx
bericht_final.docx
bericht_final_wirklich_final.docx
bericht_final_wirklich_final_Anna.docx
\end{lstlisting}

Bei einem einzelnen Dokument ist das lästig. Bei Software mit Hunderten von Dateien und
mehreren Beteiligten wird es unbeherrschbar. Vier Fragen lassen sich ohne Werkzeug kaum
noch beantworten:

\begin{enumerate}
  \item \textbf{Was hat sich geändert?} Welche Zeilen sind zwischen gestern und heute anders?
  \item \textbf{Wer hat es geändert und warum?} Wer hat die Zeile eingefügt, die seit Montag
        einen Fehler verursacht, und was war die Absicht?
  \item \textbf{Wie komme ich zurück?} Gestern lief das Programm noch. Wie stelle ich den
        damaligen Zustand wieder her, ohne die Arbeit von heute zu verlieren?
  \item \textbf{Wie arbeiten mehrere Personen gleichzeitig?} Anna und Ben ändern am selben
        Tag dieselbe Datei. Wessen Version gilt?
\end{enumerate}

\section{Die Lösung: ein Versionsverwaltungssystem}

\begin{definition}{Versionsverwaltung}
Ein \begriff{Versionsverwaltungssystem} (englisch \emph{Version Control System}, VCS)
speichert die Geschichte eines Projekts als Folge von Zwischenständen. Zu jedem
Zwischenstand werden Autor, Zeitpunkt und eine Beschreibung festgehalten. Jeder frühere
Stand lässt sich jederzeit wiederherstellen und mit anderen vergleichen.
\end{definition}

Damit sind die vier Fragen beantwortet: Das System weiss, \emph{was} sich geändert hat,
\emph{wer} es war und \emph{warum}, es kann zu jedem Stand \emph{zurück}, und es kann die
Arbeit mehrerer Personen \emph{zusammenführen}.

Ein Versionsverwaltungssystem ersetzt dabei kein Backup und keinen Cloud-Speicher wie
OneDrive. Der Unterschied liegt in der Absicht: Ein Cloud-Speicher synchronisiert
\emph{automatisch} jede Änderung. Bei der Versionsverwaltung entscheidet der Mensch
\emph{bewusst}, wann ein Stand gut genug ist, um festgehalten zu werden, und beschreibt ihn.
Dadurch entsteht eine lesbare Geschichte statt eines Stroms von Dateiversionen.

\section{Git}

\begriff{Git} ist das mit Abstand verbreitetste Versionsverwaltungssystem. Es ist kostenlos,
quelloffen und läuft auf jedem Betriebssystem.

\begin{hintergrund}
Linus Torvalds, der Erfinder von Linux, entwickelte Git im Jahr 2005 in wenigen Wochen.
Das Linux-Projekt hatte damals Tausende Mitwirkende auf der ganzen Welt und verlor über
Nacht das Recht, sein bisheriges Werkzeug kostenlos zu nutzen. Torvalds' Anforderungen waren
entsprechend hoch: Git sollte sehr schnell sein, auch ohne Internetverbindung funktionieren,
und kein einzelner Server sollte das Projekt lahmlegen können. Diese Entscheidungen prägen
Git bis heute.
\end{hintergrund}

Die wichtigste dieser Entscheidungen: Git ist \begriff{verteilt}. Ältere Systeme hatten
einen zentralen Server, auf dem die Geschichte lag. Ohne Verbindung zum Server konnte man
nicht arbeiten. Bei Git besitzt dagegen \textbf{jede Kopie die vollständige Geschichte}.
Man kann offline Zwischenstände speichern, die Geschichte durchsuchen und Versionen
vergleichen. Synchronisiert wird nur, wenn man es will.

\section{Git und GitHub: zwei verschiedene Dinge}

Einsteiger verwechseln die beiden häufig, dabei ist die Trennung einfach:

\begin{center}
\renewcommand{\arraystretch}{1.35}
\begin{tabularx}{\linewidth}{@{}l>{\raggedright\arraybackslash}X>{\raggedright\arraybackslash}X@{}}
  \toprule
                     & \textbf{Git} & \textbf{GitHub} \\
  \midrule
  Was es ist         & ein Programm auf dem eigenen Computer & ein Internetdienst (Website) \\
  Aufgabe            & Geschichte speichern, vergleichen, zusammenführen & Kopien im Internet aufbewahren und teilen \\
  Zusätzliches       & –  & Pull Requests, Code Review, automatische Tests, Fehlerverwaltung \\
  Gehört zu          & freie Software, keiner Firma & Microsoft (seit 2018) \\
  Alternativen       & –  & GitLab, Bitbucket, eigener Server \\
  \bottomrule
\end{tabularx}
\end{center}

Ein Vergleich hilft: Git ist wie die Fotokamera, GitHub wie das Online-Fotoalbum, in dem
man Bilder sichert und mit anderen teilt. Die Kamera funktioniert auch ohne Album. Das Album
ist aber der Ort, an dem die Zusammenarbeit stattfindet.

\begin{praxis}
In Unternehmen ist die Plattform (GitHub, GitLab oder Bitbucket) das Zentrum der
Entwicklungsarbeit. Dort wird nicht nur Code aufbewahrt: Jede Änderung wird dort von
Kolleginnen und Kollegen geprüft (\emph{Code Review}), automatisch getestet und erst dann
übernommen. Diese Abläufe werden in Kapitel~\ref{kap:pr} behandelt.
\end{praxis}

\begin{zusammenfassung}
\begin{itemize}[nosep]
  \item Versionsverwaltung beantwortet: Was hat sich geändert, wer war es und warum,
        wie komme ich zurück, wie arbeiten mehrere Personen zusammen?
  \item Git ist das Werkzeug auf dem eigenen Computer, GitHub der Dienst im Internet.
  \item Git ist verteilt: Jede Kopie enthält die ganze Geschichte.
\end{itemize}
\end{zusammenfassung}

% ---------------------------------------------------------------------
\chapter{Wie Git denkt: das mentale Modell}
\label{kap:modell}

\begin{lernziele}
\begin{itemize}[nosep]
  \item die vier Orte kennen, zwischen denen Änderungen wandern,
  \item verstehen, was ein Commit technisch ist,
  \item wissen, was Branches, HEAD und Remotes sind,
  \item jeden Git-Befehl einem Weg zwischen den vier Orten zuordnen können.
\end{itemize}
\end{lernziele}

Die meisten Schwierigkeiten mit Git entstehen nicht durch vergessene Befehle, sondern durch
ein unklares Bild davon, \emph{wo} sich eine Änderung gerade befindet. Dieses Kapitel
vermittelt dieses Bild. Es ist das wichtigste Kapitel des Handbuchs.

\section{Die vier Orte}

Eine Änderung durchläuft auf ihrem Weg zu GitHub vier Orte. Jeder Git-Befehl bewegt
Änderungen von einem Ort zum nächsten oder zeigt an, was wo liegt.

\begin{center}
\begin{tikzpicture}[
    ort/.style={draw=hsluBlau, fill=hsluBlau!6, rounded corners=3pt, thick,
                minimum width=3.0cm, minimum height=1.5cm, align=center, font=\small},
    weg/.style={-{Stealth[length=2.5mm]}, thick},
    lab/.style={font=\ttfamily\footnotesize, fill=white, inner sep=1pt}
  ]
  \node[ort] (arbeit) at (0,0)    {\textbf{Arbeitsordner}\\deine Dateien};
  \node[ort] (vor)    at (3.9,0)  {\textbf{Staging Area}\\Vormerkung};
  \node[ort] (repo)   at (7.8,0)  {\textbf{lokales Repo}\\Geschichte auf\\deinem PC};
  \node[ort] (gh)     at (11.7,0) {\textbf{Remote}\\GitHub};

  \draw[weg] ([yshift=4mm]arbeit.east) -- node[lab, above=1mm] {add}     ([yshift=4mm]vor.west);
  \draw[weg] ([yshift=-4mm]vor.west)   -- node[lab, below=1mm] {restore} ([yshift=-4mm]arbeit.east);
  \draw[weg] (vor.east)                -- node[lab, above=1mm] {commit}  (repo.west);
  \draw[weg] ([yshift=4mm]repo.east)   -- node[lab, above=1mm] {push}    ([yshift=4mm]gh.west);
  \draw[weg] ([yshift=-4mm]gh.west)    -- node[lab, below=1mm] {pull}    ([yshift=-4mm]repo.east);
\end{tikzpicture}
\end{center}

\begin{description}[style=nextline, leftmargin=1.2cm]
  \item[1. Arbeitsordner \normalfont(\emph{working tree})]
    Die Dateien, die im Explorer und im Editor sichtbar sind. Hier wird gearbeitet. Git
    beobachtet diesen Ordner, verändert aber nichts von sich aus.
  \item[2. Staging Area \normalfont(Vormerkung, auch \emph{index})]
    Eine Zwischenablage für den nächsten Zwischenstand. Man legt dort gezielt die
    Änderungen ab, die zusammengehören.
  \item[3. Lokales Repository \normalfont(\emph{repository}, kurz \emph{Repo})]
    Die gespeicherte Geschichte des Projekts. Sie liegt im versteckten Ordner \cmd{.git}
    innerhalb des Projektordners. Wer diesen Ordner löscht, löscht die Geschichte.
  \item[4. Remote \normalfont(entfernte Kopie)]
    Eine Kopie des Repositorys auf einem anderen Rechner, meist auf GitHub. Sie dient als
    Sicherung und als gemeinsamer Treffpunkt für alle Beteiligten.
\end{description}

\subsection*{Warum gibt es die Staging Area?}

Die Staging Area irritiert Einsteiger am meisten: Warum nicht direkt speichern? Der Grund
ist, dass bei der Arbeit oft mehrere Dinge gleichzeitig entstehen. Angenommen, an einem
Abend wurde eine Übungsaufgabe gelöst \emph{und} nebenbei ein Tippfehler in der
Dokumentation korrigiert. Das sind zwei verschiedene Anliegen. Mit der Staging Area lassen
sie sich als zwei getrennte, sauber beschriebene Zwischenstände festhalten.

Ein Vergleich: Die Staging Area ist wie ein Einkaufskorb. Man legt hinein, was zusammengehört,
und erst an der Kasse, beim \emph{Commit}, wird verbucht.

\section{Der Commit: ein Foto des Projekts}

\begin{definition}{Commit}
Ein \begriff{Commit} ist ein gespeicherter Zwischenstand des gesamten Projekts. Er enthält
\begin{itemize}[nosep]
  \item den Inhalt aller Dateien zu diesem Zeitpunkt,
  \item Autor und Zeitpunkt,
  \item eine Beschreibung, die \emph{Commit-Nachricht},
  \item einen Verweis auf den vorherigen Commit.
\end{itemize}
\end{definition}

Ein Commit ist also kein „Änderungsprotokoll“, sondern ein vollständiges \textbf{Foto} des
Projekts. Weil jeder Commit auf seinen Vorgänger verweist, entsteht eine Kette, die
\begriff{Geschichte} (\emph{history}):

\begin{center}
\begin{tikzpicture}[
    commit/.style={circle, draw=hsluBlau, fill=hsluBlau!10, thick, minimum size=9mm,
                   font=\ttfamily\scriptsize},
    vor/.style={-{Stealth[length=2mm]}, thick, gray}
  ]
  \node[commit] (c1) at (0,0)   {98bc};
  \node[commit] (c2) at (2.2,0) {19ce};
  \node[commit] (c3) at (4.4,0) {3b1a};
  \node[commit] (c4) at (6.6,0) {7f02};
  \draw[vor] (c2) -- (c1);
  \draw[vor] (c3) -- (c2);
  \draw[vor] (c4) -- (c3);
  \node[font=\small, below=2mm of c1] {ältester};
  \node[font=\small, below=2mm of c4] {neuester};
\end{tikzpicture}
\end{center}

Die Pfeile zeigen rückwärts: Jeder Commit kennt seinen Vorgänger, nicht seinen Nachfolger.

\subsection*{Der Hash: die Adresse eines Commits}

Jeder Commit erhält eine eindeutige Kennung, den \begriff{Hash}, zum Beispiel
\cmd{3b1a03a5c2e0...}. Er besteht aus 40 Zeichen. Im Alltag genügen die ersten sieben, etwa
\cmd{3b1a03a}. Mit dem Hash lässt sich jeder Commit genau ansprechen, zum Beispiel um ihn
anzuzeigen oder rückgängig zu machen.

\begin{hintergrund}
Der Hash wird aus dem Inhalt des Commits berechnet, also aus Dateien, Autor, Zeitpunkt,
Nachricht \emph{und dem Hash des Vorgängers}. Ändert sich auch nur ein Zeichen irgendwo in
der Geschichte, ändern sich alle nachfolgenden Hashes. Darum ist die Geschichte in Git
praktisch fälschungssicher, und darum lassen sich bereits geteilte Commits nicht einfach
„umschreiben“, ohne dass alle anderen es merken (siehe Kapitel~\ref{kap:fehler}).
\end{hintergrund}

\section{Branches: parallele Arbeitslinien}

\begin{definition}{Branch}
Ein \begriff{Branch} (deutsch: Zweig) ist eine eigene Arbeitslinie in der Geschichte.
Technisch ist er nur ein beweglicher Zeiger auf einen Commit. Mit jedem neuen Commit rückt
der Zeiger des aktuellen Branches automatisch nach vorn.
\end{definition}

Jedes Repository hat einen Hauptbranch, heute meist \cmd{main} genannt (früher
\cmd{master}). Er enthält den gültigen, funktionierenden Stand. Neue Arbeit geschieht auf
einem eigenen Branch, der von \cmd{main} abzweigt:

\begin{center}
\begin{tikzpicture}[
    commit/.style={circle, draw=hsluBlau, fill=hsluBlau!10, thick, minimum size=7mm},
    neu/.style={circle, draw=bspFarbe, fill=bspFarbe!15, thick, minimum size=7mm},
    vor/.style={-{Stealth[length=2mm]}, thick, gray},
    zeiger/.style={rounded corners=2pt, font=\ttfamily\small, inner sep=3pt, text=white}
  ]
  \node[commit] (c1) at (0,0) {};
  \node[commit] (c2) at (1.8,0) {};
  \node[commit] (c3) at (3.6,0) {};
  \node[neu]    (d1) at (5.4,1.2) {};
  \node[neu]    (d2) at (7.2,1.2) {};
  \draw[vor] (c2) -- (c1);
  \draw[vor] (c3) -- (c2);
  \draw[vor] (d1) -- (c3);
  \draw[vor] (d2) -- (d1);
  \node[zeiger, fill=hsluBlau, below=4mm of c3] (m) {main};
  \draw[-{Stealth}, hsluBlau, thick] (m) -- (c3);
  \node[zeiger, fill=bspFarbe, above=4mm of d2] (a) {ads};
  \draw[-{Stealth}, bspFarbe, thick] (a) -- (d2);
  \node[font=\small\ttfamily, right=3mm of a, text=gray] (h) {HEAD};
  \draw[-{Stealth}, gray, thick] (h) -- (a);
\end{tikzpicture}
\end{center}

Hier zeigt \cmd{main} weiterhin auf den letzten gültigen Stand. Auf dem Branch \cmd{ads}
sind zwei neue Commits entstanden. \cmd{main} bleibt davon unberührt, bis die Arbeit
geprüft und übernommen wird. Dieses Übernehmen heisst \begriff{Merge} (zusammenführen).

\begin{definition}{HEAD}
\begriff{HEAD} ist der Zeiger auf „hier stehe ich gerade“. Er zeigt auf den aktuellen
Branch. Wechselt man den Branch, wandert HEAD mit, und Git passt die Dateien im
Arbeitsordner an den Stand dieses Branches an.
\end{definition}

\begin{praxis}
Weil ein Branch nur ein Zeiger ist, kostet er praktisch nichts: Anlegen, Wechseln und
Löschen dauert Sekundenbruchteile. Darum arbeiten Teams mit sehr vielen kurzlebigen
Branches, typischerweise einem pro Aufgabe oder Fehlerbehebung. Ein Branch lebt oft nur
wenige Stunden oder Tage, bis er in \cmd{main} übernommen und wieder gelöscht wird.
\end{praxis}

\section{Remotes: die Kopie im Internet}

Ein \begriff{Remote} ist ein anderes Repository, mit dem das eigene synchronisiert wird.
Beim Klonen von GitHub erhält es automatisch den Namen \cmd{origin} (Ursprung). Mit zwei
Befehlen wird abgeglichen:

\begin{itemize}
  \item \cmd{git push} schickt eigene Commits zum Remote.
  \item \cmd{git pull} holt fremde Commits vom Remote.
\end{itemize}

Git merkt sich zusätzlich, auf welchem Stand das Remote beim letzten Abgleich war. Dieser
gespeicherte Stand heisst zum Beispiel \cmd{origin/main}. Daher stammen Meldungen wie
\emph{„Your branch is ahead of \textquotesingle{}origin/main\textquotesingle{} by 1 commit“}:
Der eigene \cmd{main} hat einen Commit mehr als der zuletzt bekannte Stand auf GitHub.

\begin{zusammenfassung}
\begin{itemize}[nosep]
  \item Änderungen wandern: Arbeitsordner $\rightarrow$ Staging Area $\rightarrow$
        lokales Repo $\rightarrow$ Remote.
  \item Ein Commit ist ein vollständiges Foto des Projekts mit Autor, Zeit, Nachricht und
        Verweis auf den Vorgänger. Seine Adresse ist der Hash.
  \item Ein Branch ist ein beweglicher Zeiger auf einen Commit. \cmd{main} ist der gültige Stand.
  \item HEAD zeigt, wo man gerade steht. \cmd{origin} ist die Kopie auf GitHub.
\end{itemize}
\end{zusammenfassung}

% ---------------------------------------------------------------------
\chapter{Befehle lesen und verstehen}
\label{kap:lesen}

\begin{lernziele}
\begin{itemize}[nosep]
  \item wissen, was ein Terminal ist und was der aktuelle Ordner bedeutet,
  \item jeden Befehl in seine Bestandteile zerlegen können,
  \item den Unterschied zwischen \cmd{-x} und \cmd{--wort} begründen können,
  \item sich zu jedem Befehl selbst Hilfe holen können.
\end{itemize}
\end{lernziele}

\section{Das Terminal}

Git wird meist über ein \begriff{Terminal} bedient, ein Fenster, in das man Befehle als Text
eingibt. Unter Windows sind das zum Beispiel \emph{PowerShell} oder \emph{Git Bash}, in VS Code
das Terminal am unteren Rand. Grafische Oberflächen für Git gibt es auch, sie bauen aber auf
denselben Befehlen auf. Wer die Befehle versteht, versteht auch jede Oberfläche.

Das Terminal zeigt vor dem Cursor eine Zeile, die \begriff{Eingabeaufforderung}
(\emph{prompt}). Sie verrät, wo man sich befindet:

\begin{lstlisting}[style=ausgabe]
PS D:\github\hslu-bachelor-ai-ml\2026-FA\ADS>
bahram@PC MINGW64 /d/github/hslu-bachelor-ai-ml (main)
\end{lstlisting}

Die erste Zeile stammt aus PowerShell, die zweite aus Git Bash. Git Bash zeigt in Klammern
zusätzlich den aktuellen Branch an, hier \cmd{main}.

\subsection*{Der aktuelle Ordner}

Jedes Terminal „steht“ in einem Ordner, dem \begriff{aktuellen Arbeitsverzeichnis}. Alle
Pfade, die man angibt, gelten \textbf{relativ zu diesem Ordner}. Das ist der häufigste Grund,
warum ein Befehl „nicht funktioniert“: Man steht im falschen Ordner.

\begin{lstlisting}[style=befehle]
cd D:\github\hslu-bachelor-ai-ml       # in den Ordner wechseln (change directory)
cd 2026-FA\ADS                         # relativ: einen Unterordner tiefer
cd ..                                  # einen Ordner nach oben
pwd                                    # zeigen, wo man steht (print working dir.)
\end{lstlisting}

Git-Befehle wirken immer auf das Repository, in dem der aktuelle Ordner liegt. Git sucht
dazu vom aktuellen Ordner aus nach oben, bis es den Ordner \cmd{.git} findet. Deshalb
funktionieren Git-Befehle in jedem Unterordner eines Repositorys.

\section{Aufbau eines Befehls}

Alle Git-Befehle folgen demselben Muster:

\begin{center}
\begin{tikzpicture}[
    teil/.style={font=\ttfamily\large, inner sep=2pt},
    erkl/.style={font=\small, align=center},
    klammer/.style={decorate, decoration={brace, amplitude=4pt, mirror}, thick, hsluBlau}
  ]
  \node[teil] (p)  at (0,0)    {git};
  \node[teil] (u)  at (2.2,0)  {diff};
  \node[teil] (o1) at (4.6,0)  {--cached};
  \node[teil] (o2) at (7.0,0)  {--stat};
  \node[teil] (t)  at (9.0,0)  {--};
  \node[teil] (a)  at (11.2,0) {datei.py};
  \draw[klammer] (p.south west)  -- (p.south east)  node[midway, below=6pt, erkl] {Programm};
  \draw[klammer] (u.south west)  -- (u.south east)  node[midway, below=6pt, erkl] {Unterbefehl:\\was getan wird};
  \draw[klammer] (o1.south west) -- (o2.south east) node[midway, below=6pt, erkl] {Optionen:\\wie es getan wird};
  \draw[klammer] (t.south west)  -- (t.south east)  node[midway, below=6pt, erkl] {Trenner};
  \draw[klammer] (a.south west)  -- (a.south east)  node[midway, below=6pt, erkl] {Argument:\\worauf};
\end{tikzpicture}
\end{center}

\begin{description}[style=nextline, leftmargin=1.2cm]
  \item[Programm] \cmd{git} oder \cmd{gh}. Das Terminal startet dieses Programm.
  \item[Unterbefehl] die Tätigkeit: \cmd{status}, \cmd{add}, \cmd{commit}, \ldots{} Git hat
    über hundert Unterbefehle. Für den Alltag genügen etwa fünfzehn.
  \item[Optionen] verändern \emph{wie} der Unterbefehl arbeitet. Sie beginnen mit einem
    oder zwei Strichen. Die meisten sind freiwillig.
  \item[Argumente] sagen, \emph{worauf} sich der Befehl bezieht, zum Beispiel Dateien,
    Ordner oder Branch-Namen.
\end{description}

\section{Ein Strich oder zwei Striche}

Die Schreibweise der Optionen ist keine Eigenheit von Git, sondern eine Konvention, die
fast alle Kommandozeilen-Programme seit den 1990er-Jahren teilen:

\begin{center}
\renewcommand{\arraystretch}{1.3}
\begin{tabularx}{\linewidth}{@{}l>{\ttfamily}l>{\raggedright\arraybackslash}X@{}}
  \toprule
  \textbf{Form} & \normalfont\textbf{Beispiel} & \textbf{Zweck} \\
  \midrule
  ein Strich, ein Buchstabe & -m, -s, -b & Kurzform: schnell zu tippen, für den Alltag \\
  zwei Striche, ganzes Wort & --message, --short & Langform: selbsterklärend, für Skripte und Dokumentation \\
  \bottomrule
\end{tabularx}
\end{center}

Viele Optionen gibt es in beiden Formen: \cmd{-m} und \cmd{--message} bedeuten dasselbe.

Die zwei Striche haben einen technischen Grund. Kurzformen dürfen \textbf{zusammengezogen}
werden: \cmd{-sb} bedeutet dasselbe wie \cmd{-s -b}. Würde man ganze Wörter mit nur einem
Strich schreiben, wäre \cmd{-stat} mehrdeutig: die Option \cmd{stat}, oder die vier
Optionen \cmd{-s -t -a -t}? Der doppelte Strich beseitigt diese Mehrdeutigkeit.

\subsection*{Optionen mit Wert}

Manche Optionen brauchen einen Wert. Bei der Kurzform folgt er nach einem Leerzeichen, bei
der Langform meist nach einem Gleichheitszeichen:

\begin{lstlisting}[style=befehle]
git commit -m "Stack implementiert"     # Kurzform
git log --format="%h %s"                # Langform
\end{lstlisting}

Enthält ein Wert Leerzeichen, gehört er in \textbf{Anführungszeichen}. Das Terminal trennt
Eingaben an Leerzeichen auf. Ohne Anführungszeichen würde
\cmd{git commit -m Stack implementiert} nur \cmd{Stack} als Nachricht verwenden und
\cmd{implementiert} als Dateinamen deuten.

\subsection*{Zwei Striche ohne Wort}

Ein alleinstehendes \cmd{--} bedeutet: „Ab hier folgen keine Optionen mehr, nur noch Pfade.“
Das ist nötig, wenn ein Dateiname wie eine Option oder ein Branch aussieht:

\begin{lstlisting}[style=befehle]
git restore -- main.py         # "main.py" ist sicher eine Datei, kein Branch "main"
\end{lstlisting}

Im Alltag braucht man das selten. Es taucht aber oft in Anleitungen und Fehlermeldungen auf.

\section{Hilfe zu jedem Befehl}

Niemand kennt alle Optionen auswendig. Erfahrene Entwicklerinnen und Entwickler schlagen
ständig nach. Git liefert die Dokumentation gleich mit:

\begin{lstlisting}[style=befehle]
git diff -h                    # kurze Übersicht aller Optionen im Terminal
git help diff                  # ausführliches Handbuch im Browser
\end{lstlisting}

\begin{erfahrung}
Fehlermeldungen von Git sind ungewöhnlich hilfreich. Sie nennen fast immer die Ursache und
oft schon den passenden Befehl, besonders in Zeilen, die mit \cmd{hint:} beginnen. Wer sie
in Ruhe liest, statt sie wegzuklicken, löst die meisten Probleme selbst.
\end{erfahrung}

\begin{zusammenfassung}
\begin{itemize}[nosep]
  \item Pfade gelten relativ zum aktuellen Ordner des Terminals.
  \item Aufbau: Programm, Unterbefehl, Optionen, Argumente.
  \item \cmd{-x} ist die Kurzform, \cmd{--wort} die Langform. Kurzformen lassen sich zusammenziehen.
  \item Werte mit Leerzeichen gehören in Anführungszeichen. \cmd{--} allein trennt Optionen von Pfaden.
  \item \cmd{git <befehl> -h} und \cmd{git help <befehl>} liefern Hilfe.
\end{itemize}
\end{zusammenfassung}


\part{Arbeiten mit Git}
% ---------------------------------------------------------------------
\chapter{Einrichten}
\label{kap:einrichten}

\begin{lernziele}
\begin{itemize}[nosep]
  \item Git so einrichten, dass jeder Commit korrekt dem eigenen Namen zugeordnet wird,
  \item die GitHub-Anmeldung auf der Kommandozeile einrichten,
  \item ein bestehendes Projekt klonen oder ein neues anlegen.
\end{itemize}
\end{lernziele}

\section{Die eigene Identität}

Jeder Commit speichert, \emph{wer} ihn erstellt hat. Darum muss Git einmalig Name und
E-Mail-Adresse kennen:

\begin{lstlisting}[style=befehle]
git config --global user.name  "Bahram Alizadah"
git config --global user.email "name@example.com"
\end{lstlisting}

\begin{teile}
  config        & Einstellungen von Git lesen oder schreiben \\
  --global      & gilt für alle Repositorys dieses Benutzers auf diesem Computer \\
  user.name     & der Name, der in jedem Commit als Autor erscheint \\
  user.email    & die E-Mail-Adresse des Autors \\
\end{teile}

Die E-Mail-Adresse sollte dieselbe sein wie im GitHub-Konto. Nur dann kann GitHub die
Commits dem Profil zuordnen und mit dem eigenen Bild anzeigen.

\begin{praxis}
In Unternehmen ist die korrekte Identität keine Formsache. Über die Geschichte lässt sich
für jede Codezeile nachvollziehen, wer sie wann eingeführt hat. Das hilft bei der
Fehlersuche („Wen frage ich, was hier gemeint war?“) und ist in regulierten Branchen wie
Banken oder Medizintechnik oft vorgeschrieben.
\end{praxis}

\section{Zwei sinnvolle Grundeinstellungen}

\begin{lstlisting}[style=befehle]
git config --global init.defaultBranch main
git config --global core.autocrlf true
\end{lstlisting}

\begin{teile}
  init.defaultBranch main & neue Repositorys erhalten den Hauptbranch \cmd{main}, wie es heute üblich ist \\
  core.autocrlf true      & passt unter Windows die Zeilenenden automatisch an (siehe Abschnitt~\ref{sec:zeilenenden}) \\
\end{teile}

Alle aktuellen Einstellungen zeigt:

\begin{lstlisting}[style=befehle]
git config --global --list
\end{lstlisting}

\section{Anmeldung bei GitHub}

Für Pull Requests auf der Kommandozeile gibt es das Programm \cmd{gh} (GitHub CLI). Es
muss einmal pro Computer angemeldet werden:

\begin{lstlisting}[style=befehle]
gh auth login --hostname github.com --git-protocol https --web
\end{lstlisting}

\begin{teile}
  auth login              & Anmeldung starten \\
  --hostname github.com   & bei welchem Dienst (es gibt auch firmeneigene GitHub-Server) \\
  --git-protocol https    & wie Daten übertragen werden; HTTPS ist der einfachste Weg \\
  --web                   & Anmeldung im Browser mit einem einmaligen Code \\
\end{teile}

Danach zeigt \cmd{gh} einen achtstelligen Code an. Im Browser wird er auf
\cmd{github.com/login/device} eingegeben und mit \emph{Authorize} bestätigt. \cmd{gh}
erhält dadurch einen \begriff{Token}, einen digitalen Schlüssel, mit dem es im Namen des
Benutzers handeln darf. Windows speichert ihn sicher in der Anmeldeinformationsverwaltung.

\begin{fehler}
In \emph{Git Bash} bricht \cmd{gh auth login} ohne Optionen mit „could not prompt“ ab.
Git Bash kann die Auswahlmenüs von \cmd{gh} nicht darstellen. Mit den oben gezeigten
Optionen entfallen die Menüs. Alternativ funktioniert PowerShell.
\end{fehler}

\section{Ein Repository beschaffen}

Es gibt zwei Wege zu einem Repository.

\textbf{Klonen} kopiert ein bestehendes Repository von GitHub auf den eigenen Computer,
samt vollständiger Geschichte. Das ist der Normalfall:

\begin{lstlisting}[style=befehle]
cd D:\github
gh repo clone gladtobethere/hslu-bachelor-ai-ml
\end{lstlisting}

Es entsteht der Ordner \texttt{D:\textbackslash{}github\textbackslash{}hslu-bachelor-ai-ml}, und \cmd{origin} zeigt
automatisch auf GitHub. Mit reinem Git lautet der Befehl \cmd{git clone <URL>}.

\textbf{Neu anlegen} macht einen bestehenden Ordner zu einem Repository:

\begin{lstlisting}[style=befehle]
git init                       # legt den Ordner .git an
\end{lstlisting}

\begin{erfahrung}
Ein Ordner pro Repository, an einem festen Ort, zum Beispiel \texttt{D:\textbackslash{}github\textbackslash{}<repo>}.
Mehrere Kopien desselben Repositorys auf einem Computer führen früher oder später dazu,
dass man in der falschen Kopie arbeitet und Änderungen „verschwinden“.
\end{erfahrung}

\begin{zusammenfassung}
\begin{itemize}[nosep]
  \item Einmalig \cmd{user.name} und \cmd{user.email} setzen, passend zum GitHub-Konto.
  \item \cmd{gh auth login} meldet die Kommandozeile bei GitHub an.
  \item Klonen holt ein bestehendes Repository, \cmd{git init} legt ein neues an.
\end{itemize}
\end{zusammenfassung}

% ---------------------------------------------------------------------
\chapter{Der Arbeitsablauf}
\label{kap:ablauf}

\begin{lernziele}
\begin{itemize}[nosep]
  \item den Feature-Branch-Workflow kennen, wie ihn die meisten Teams verwenden,
  \item verstehen, \emph{warum} nie direkt auf \cmd{main} gearbeitet wird,
  \item einen vollständigen Arbeitszyklus von Anfang bis Ende durchführen können.
\end{itemize}
\end{lernziele}

\section{Warum nicht einfach auf main arbeiten?}

Technisch kann man auf \cmd{main} committen und direkt pushen. Für ein Wegwerfprojekt geht
das. In der Industrie ist es dagegen fast überall verboten, und zwar aus drei Gründen:

\begin{enumerate}
  \item \textbf{main muss immer funktionieren.} Aus \cmd{main} wird die Software gebaut,
        die Kundinnen und Kunden verwenden. Ein halbfertiger Zwischenstand darf dort nie landen.
  \item \textbf{Vier-Augen-Prinzip.} Jede Änderung wird vor der Übernahme von einer zweiten
        Person geprüft. Das findet Fehler, verbreitet Wissen im Team und verhindert, dass
        jemand allein unbemerkt etwas Falsches einbaut.
  \item \textbf{Automatische Prüfung.} Vor der Übernahme laufen automatisch Tests. Schlägt
        einer fehl, wird die Änderung nicht übernommen.
\end{enumerate}

Diese Prüfungen brauchen einen Ort, an dem die Änderung \emph{wartet}, bevor sie in
\cmd{main} kommt: den eigenen Branch und den Pull Request.

\begin{praxis}
Auf GitHub lässt sich \cmd{main} technisch sperren (\emph{branch protection}). Ein direkter
Push wird dann abgewiesen, und eine Übernahme ist nur über einen Pull Request mit
Freigabe und grünen Tests möglich. In professionellen Teams ist das die Regel. Wer sich
den Ablauf schon bei eigenen Projekten angewöhnt, arbeitet vom ersten Tag an wie im Team.
\end{praxis}

\section{Der Feature-Branch-Workflow}

Der verbreitetste Ablauf heisst \begriff{Feature-Branch-Workflow}: Jede Aufgabe, ob neue
Funktion, Fehlerbehebung oder Übung, erhält einen eigenen kurzlebigen Branch.

\begin{center}
\begin{tikzpicture}[
    schritt/.style={draw=hsluBlau, fill=hsluBlau!6, rounded corners=3pt, thick,
                    minimum width=2.45cm, minimum height=1.05cm, align=center, font=\small},
    weg/.style={-{Stealth[length=2.2mm]}, thick}
  ]
  \node[schritt] (s1) at (0,0)    {aktuellen\\Stand holen};
  \node[schritt] (s2) at (3.0,0)  {Branch\\anlegen};
  \node[schritt] (s3) at (6.0,0)  {arbeiten,\\committen};
  \node[schritt] (s4) at (9.0,0)  {pushen};
  \node[schritt] (s5) at (9.0,-1.8) {Pull Request\\öffnen};
  \node[schritt] (s6) at (6.0,-1.8) {prüfen,\\testen};
  \node[schritt] (s7) at (3.0,-1.8) {in main\\übernehmen};
  \node[schritt] (s8) at (0,-1.8)  {Branch\\löschen};
  \draw[weg] (s1) -- (s2);  \draw[weg] (s2) -- (s3);  \draw[weg] (s3) -- (s4);
  \draw[weg] (s4) -- (s5);  \draw[weg] (s5) -- (s6);  \draw[weg] (s6) -- (s7);
  \draw[weg] (s7) -- (s8);  \draw[weg, dashed, gray] (s8) -- (s1);
\end{tikzpicture}
\end{center}

\section{Ein vollständiger Arbeitszyklus}

Das folgende Beispiel zeigt einen Abend, an dem eine Übung im Modul ADS gelöst wird. Jeder
Befehl wird in Kapitel~\ref{kap:befehle} im Detail erklärt.

\subsection*{1. Aktuellen Stand holen}

\begin{lstlisting}[style=befehle]
git switch main
git pull
\end{lstlisting}

Wer auf einem veralteten Stand beginnt, riskiert später Konflikte mit Änderungen, die in
der Zwischenzeit entstanden sind, zum Beispiel auf einem anderen Gerät.

\subsection*{2. Branch anlegen}

\begin{lstlisting}[style=befehle]
git switch -c ads
\end{lstlisting}

Ab jetzt landen alle Commits auf \cmd{ads}. \cmd{main} bleibt unberührt.

\subsection*{3. Arbeiten und committen}

Dateien bearbeiten, speichern, dann:

\begin{lstlisting}[style=befehle]
git status                             # was hat sich geändert?
git add python                         # Änderungen auswählen
git diff --cached --stat               # Auswahl kontrollieren
git commit -m "ADS U02: Stack mit verketteter Liste implementiert"
\end{lstlisting}

Dieser Schritt darf sich mehrmals wiederholen. Mehrere kleine Commits sind besser als ein
grosser (siehe Kapitel~\ref{kap:commits}).

\subsection*{4. Pushen}

\begin{lstlisting}[style=befehle]
git push -u origin ads
\end{lstlisting}

Jetzt liegt die Arbeit auch auf GitHub und ist gegen den Ausfall des eigenen Computers gesichert.

\subsection*{5. Pull Request öffnen, prüfen, übernehmen}

\begin{lstlisting}[style=befehle]
gh pr create --fill
gh pr merge --merge --delete-branch
\end{lstlisting}

Im Team würde zwischen diesen beiden Befehlen eine Kollegin den Pull Request prüfen und
freigeben. Bei eigenen Projekten übernimmt man diese Rolle selbst und schaut sich die
Änderungen auf GitHub noch einmal in Ruhe an.

\subsection*{6. Abschliessen}

\begin{lstlisting}[style=befehle]
git pull
\end{lstlisting}

\cmd{gh pr merge --delete-branch} hat bereits auf \cmd{main} zurückgewechselt. \cmd{git pull}
holt den neuen Stand von \cmd{main}, der jetzt die eigene Arbeit enthält. Der Kreis ist
geschlossen.

\begin{erfahrung}
\textbf{Merkhilfe für das Ende:} \emph{push} (hochladen) $\rightarrow$ \emph{create}
(Antrag stellen) $\rightarrow$ \emph{merge} (Antrag annehmen) $\rightarrow$ \emph{pull}
(Ergebnis holen). Wer \emph{create} vergisst, erhält beim Mergen die Meldung
„no pull requests found“: Ohne Antrag gibt es nichts anzunehmen.
\end{erfahrung}

\section{Branch-Namen}

In Teams folgen Branch-Namen oft einem Schema, damit man auf einen Blick sieht, worum es
geht, zum Beispiel \cmd{feature/login}, \cmd{fix/absturz-beim-speichern} oder
\cmd{JIRA-1234-export}. Für eigene Projekte genügt ein kurzes Wort ohne Leerzeichen und
Umlaute: \cmd{ads}, \cmd{lial}, \cmd{latex}. Da der Branch nach dem Merge gelöscht wird,
kann derselbe Name immer wieder verwendet werden.

\begin{zusammenfassung}
\begin{itemize}[nosep]
  \item \cmd{main} enthält immer einen funktionierenden Stand und wird nie direkt bearbeitet.
  \item Jede Aufgabe erhält einen eigenen Branch, der per Pull Request übernommen wird.
  \item Zyklus: \cmd{pull} $\rightarrow$ \cmd{switch -c} $\rightarrow$ \cmd{add}/\cmd{commit}
        $\rightarrow$ \cmd{push} $\rightarrow$ \cmd{pr create} $\rightarrow$ \cmd{pr merge}
        $\rightarrow$ \cmd{pull}.
\end{itemize}
\end{zusammenfassung}

% ---------------------------------------------------------------------
\chapter{Die Befehle im Einzelnen}
\label{kap:befehle}

\begin{lernziele}
\begin{itemize}[nosep]
  \item für jeden Alltagsbefehl wissen, welchen Weg im mentalen Modell er geht,
  \item jede Option eines Befehls einzeln verstehen und gezielt kombinieren,
  \item die Ausgaben von \cmd{status}, \cmd{diff}, \cmd{log} und \cmd{push} lesen können.
\end{itemize}
\end{lernziele}

Die Befehle sind in der Reihenfolge des Arbeitszyklus geordnet. Jeder Abschnitt beginnt mit
dem Zweck, zerlegt dann den Befehl in seine Teile und schliesst mit Hinweisen aus der Praxis.

% =====================================================================
\befehl{git status}{den Zustand abfragen}

\textbf{Zweck:} zeigt, welche Änderungen sich in welchem der vier Orte befinden. Der Befehl
verändert nichts. Er ist der meistverwendete Git-Befehl überhaupt.

\begin{lstlisting}[style=befehle]
git status
\end{lstlisting}

\begin{lstlisting}[style=ausgabe]
On branch main
Your branch is up to date with 'origin/main'.

Changes to be committed:
        modified:   python/uebung01/al/aufgabe03/recursive_sum.py

Changes not staged for commit:
        modified:   README.md

Untracked files:
        notizen.txt
\end{lstlisting}

Die Ausgabe gliedert sich nach den Orten des mentalen Modells:

\begin{center}
\renewcommand{\arraystretch}{1.3}
\begin{tabularx}{\linewidth}{@{}>{\raggedright\arraybackslash}p{4.6cm}>{\raggedright\arraybackslash}X@{}}
  \toprule
  \textbf{Abschnitt} & \textbf{Bedeutung} \\
  \midrule
  Changes to be committed      & in der Staging Area, kommt in den nächsten Commit (grün) \\
  Changes not staged for commit& im Arbeitsordner geändert, aber nicht vorgemerkt (rot) \\
  Untracked files              & neue Dateien, die Git noch nie gesehen hat (rot) \\
  up to date with origin/main  & gleicher Stand wie beim letzten Abgleich mit GitHub \\
  \bottomrule
\end{tabularx}
\end{center}

\subsection*{Kurzform}

\begin{lstlisting}[style=befehle]
git status -sb
\end{lstlisting}

\begin{teile}
  -s  & \textbf{s}hort: eine Zeile pro Datei statt ganzer Sätze \\
  -b  & \textbf{b}ranch: zeigt trotzdem die Zeile mit dem Branch \\
\end{teile}

\begin{lstlisting}[style=ausgabe]
## main...origin/main [ahead 1]
M  python/uebung01/al/aufgabe03/recursive_sum.py
 M README.md
?? notizen.txt
\end{lstlisting}

Die zwei Zeichen vor jedem Dateinamen sind zwei Spalten: Die \textbf{linke} beschreibt die
Staging Area, die \textbf{rechte} den Arbeitsordner. \cmd{[ahead 1]} bedeutet: ein Commit
mehr als auf GitHub.

\begin{center}
\renewcommand{\arraystretch}{1.25}
\begin{tabularx}{\linewidth}{@{}>{\ttfamily}l>{\raggedright\arraybackslash}X@{}}
  \toprule
  \normalfont\textbf{Zeichen} & \textbf{Bedeutung} \\
  \midrule
  M\textvisiblespace  & geändert und vorgemerkt \\
  \textvisiblespace M & geändert, nicht vorgemerkt \\
  MM                  & vorgemerkt und danach erneut geändert \\
  A\textvisiblespace  & neue Datei, vorgemerkt (\emph{added}) \\
  D\textvisiblespace  & gelöscht, vorgemerkt (\emph{deleted}) \\
  ??                  & unbekannte Datei (\emph{untracked}) \\
  \bottomrule
\end{tabularx}
\end{center}

\begin{erfahrung}
Vor jedem Commit \cmd{git status}. Die zwei Sekunden verhindern fast jeden Fall von „das
hätte nicht mit sollen“.
\end{erfahrung}

% =====================================================================
\befehl{git add}{Änderungen vormerken}

\textbf{Zweck:} legt Änderungen aus dem Arbeitsordner in die Staging Area.
\textbf{Weg:} Arbeitsordner $\rightarrow$ Staging Area.

\begin{lstlisting}[style=befehle]
git add python java
\end{lstlisting}

Die Argumente sind Dateien oder Ordner, \textbf{relativ zum aktuellen Ordner}. Bei einem
Ordner werden alle Änderungen darin und in allen Unterordnern vorgemerkt.

\begin{center}
\renewcommand{\arraystretch}{1.25}
\begin{tabularx}{\linewidth}{@{}>{\ttfamily}l>{\raggedright\arraybackslash}X@{}}
  \toprule
  \normalfont\textbf{Aufruf} & \textbf{Was vorgemerkt wird} \\
  \midrule
  git add datei.py & genau diese Datei \\
  git add python   & alles unter \cmd{python/} \\
  git add .        & alles im aktuellen Ordner und darunter; der Punkt bedeutet „hier“ \\
  git add -A       & alle Änderungen im ganzen Repository (\emph{all}), auch gelöschte Dateien \\
  git add -p       & interaktiv Stück für Stück (\emph{patch}); Git fragt bei jedem Abschnitt \cmd{y} oder \cmd{n} \\
  \bottomrule
\end{tabularx}
\end{center}

Vorgemerkt werden nur \textbf{Änderungen}. Unveränderte Dateien sind bereits im letzten
Commit enthalten und werden nicht erneut gespeichert.

\begin{praxis}
Erfahrene Entwicklerinnen und Entwickler verwenden \cmd{git add -A} zurückhaltend. Sobald
mehrere Anliegen gleichzeitig offen sind, merken sie gezielt Pfade vor oder nutzen
\cmd{git add -p}, um sogar innerhalb einer Datei zu trennen. So bleibt jeder Commit auf ein
Anliegen beschränkt, was das Prüfen und spätere Rückgängigmachen erleichtert.
\end{praxis}

% =====================================================================
\befehl{git diff}{Unterschiede anzeigen}

\textbf{Zweck:} zeigt Zeile für Zeile, was sich zwischen zwei Ständen unterscheidet
(\emph{diff} von \emph{difference}). Welche zwei Stände verglichen werden, hängt von den
Optionen ab:

\begin{center}
\renewcommand{\arraystretch}{1.3}
\begin{tabularx}{\linewidth}{@{}>{\ttfamily}l>{\raggedright\arraybackslash}p{4.4cm}>{\raggedright\arraybackslash}X@{}}
  \toprule
  \normalfont\textbf{Aufruf} & \textbf{vergleicht} & \textbf{beantwortet die Frage} \\
  \midrule
  git diff          & Arbeitsordner $\leftrightarrow$ Staging Area & Was habe ich geändert, aber noch nicht vorgemerkt? \\
  git diff --cached & Staging Area $\leftrightarrow$ letzter Commit & Was kommt in den nächsten Commit? \\
  git diff HEAD     & Arbeitsordner $\leftrightarrow$ letzter Commit & Was ist insgesamt seit dem letzten Commit anders? \\
  \bottomrule
\end{tabularx}
\end{center}

\begin{hintergrund}
\cmd{--cached} heisst so, weil die Staging Area in Git intern früher \emph{cache} genannt
wurde. Da der Name verwirrend ist, gibt es seit einigen Jahren das gleichbedeutende
\cmd{--staged}.
\end{hintergrund}

\subsection*{Die Ausgabe lesen}

\begin{lstlisting}[style=ausgabe]
--- a/recursive_sum.py
+++ b/recursive_sum.py
@@ -3,4 +3,6 @@
 def recursive_sum(n):
-    pass
+    if n == 0:
+        return 0
+    return n + recursive_sum(n - 1)
\end{lstlisting}

\begin{itemize}
  \item \cmd{---} kennzeichnet die alte Fassung (a), \cmd{+++} die neue (b).
  \item \cmd{@@ -3,4 +3,6 @@} bedeutet: Der Ausschnitt beginnt in Zeile 3; vorher umfasste er
        4 Zeilen, jetzt 6.
  \item Zeilen mit \cmd{-} wurden entfernt, Zeilen mit \cmd{+} hinzugefügt. Zeilen ohne
        Zeichen sind unverändert und dienen der Orientierung.
  \item Lange Ausgaben zeigt Git seitenweise: Leertaste blättert weiter, \cmd{q} beendet.
\end{itemize}

\subsection*{Nur die Übersicht: \texorpdfstring{\cmd{--stat}}{--stat}}

\begin{lstlisting}[style=befehle]
git diff --cached --stat
\end{lstlisting}

\begin{lstlisting}[style=ausgabe]
 java/uebung01/as/aufgabe03/RecursiveSum.java  | 12 +++++++---
 python/uebung01/al/aufgabe03/recursive_sum.py |  5 +++--
 2 files changed, 13 insertions(+), 4 deletions(-)
\end{lstlisting}

\cmd{--stat} (\emph{statistics}) zeigt pro Datei nur die Anzahl geänderter Zeilen und das
Verhältnis von hinzugefügten (\cmd{+}) zu entfernten (\cmd{-}) Zeilen.

Die beiden Optionen sind voneinander unabhängig: \cmd{--cached} bestimmt, \textbf{was}
verglichen wird, \cmd{--stat}, \textbf{wie} es dargestellt wird. Am besten probiert man sie
einzeln aus:

\begin{center}
\renewcommand{\arraystretch}{1.25}
\begin{tabularx}{\linewidth}{@{}>{\ttfamily}l>{\raggedright\arraybackslash}X@{}}
  \toprule
  \normalfont\textbf{Aufruf} & \textbf{Ergebnis} \\
  \midrule
  git diff                      & Details, nur nicht Vorgemerktes \\
  git diff --stat               & Übersicht, nur nicht Vorgemerktes \\
  git diff --cached             & Details, nur Vorgemerktes \\
  git diff --cached --stat      & Übersicht, nur Vorgemerktes \\
  git diff --cached --name-only & nur die Dateinamen \\
  \bottomrule
\end{tabularx}
\end{center}

% =====================================================================
\befehl{git restore}{Änderungen zurücknehmen}

\textbf{Zweck:} holt Dateien aus der Staging Area zurück oder verwirft Änderungen.
\textbf{Weg:} Staging Area $\rightarrow$ Arbeitsordner.

\begin{lstlisting}[style=befehle]
git restore --staged datei.py          # aus der Vormerkung nehmen, Datei bleibt
git restore datei.py                   # Änderungen an der Datei verwerfen
\end{lstlisting}

\begin{teile}
  --staged & nur die Vormerkung aufheben; der Inhalt der Datei bleibt unverändert \\
  (ohne)   & die Datei im Arbeitsordner auf den letzten Commit zurücksetzen \\
\end{teile}

\begin{achtung}
\cmd{git restore datei.py} \emph{ohne} \cmd{--staged} löscht alle nicht committeten
Änderungen an dieser Datei endgültig. Es gibt keinen Papierkorb. Vorher mit
\cmd{git diff datei.py} prüfen, was verloren ginge.
\end{achtung}

% =====================================================================
\befehl{git commit}{einen Zwischenstand speichern}

\textbf{Zweck:} speichert alles Vorgemerkte als neuen Commit.
\textbf{Weg:} Staging Area $\rightarrow$ lokales Repository.

\begin{lstlisting}[style=befehle]
git commit -m "ADS U01: rekursive Summe implementiert"
\end{lstlisting}

\begin{teile}
  -m & \textbf{m}essage: die Beschreibung direkt angeben; ohne \cmd{-m} öffnet sich ein Texteditor \\
\end{teile}

Der Commit existiert danach nur \textbf{lokal}. GitHub erfährt erst beim Push davon. Wie man
gute Nachrichten schreibt, behandelt Kapitel~\ref{kap:commits}.

% =====================================================================
\befehl{git log}{die Geschichte ansehen}

\textbf{Zweck:} listet die Commits auf, den neuesten zuerst.

\begin{lstlisting}[style=befehle]
git log --oneline --graph --all -10
\end{lstlisting}

\begin{teile}
  --oneline & eine Zeile pro Commit: kurzer Hash und erste Zeile der Nachricht \\
  --graph   & zeichnet die Verzweigungen der Branches als Linien \\
  --all     & alle Branches statt nur des aktuellen \\
  -10       & nur die letzten zehn Commits \\
\end{teile}

\begin{lstlisting}[style=ausgabe]
*   3b1a03a Merge pull request #1 from gladtobethere/ads
|\
| * db753be ADS: Uebungen U01-U06
|/
* 19ce440 ADS-Uebungsordner anlegen
\end{lstlisting}

Die Grafik zeigt: Auf einem Branch entstand \cmd{db753be}, der danach mit einem
Merge-Commit (\cmd{3b1a03a}) in \cmd{main} übernommen wurde.

% =====================================================================
\befehl{git switch}{Branch wechseln oder anlegen}

\textbf{Zweck:} wechselt den aktuellen Branch, HEAD wandert mit.

\begin{lstlisting}[style=befehle]
git switch main                        # auf bestehenden Branch wechseln
git switch -c ads                      # neuen Branch anlegen und wechseln
\end{lstlisting}

\begin{teile}
  -c & \textbf{c}reate: den Branch zuerst anlegen, ausgehend vom aktuellen Commit \\
\end{teile}

Nicht committete Änderungen wandern beim Anlegen eines neuen Branches mit. Wer versehentlich
auf \cmd{main} gearbeitet hat, rettet die Arbeit deshalb einfach mit \cmd{git switch -c <name>}.

\begin{hintergrund}
\cmd{git switch} gibt es seit 2019. Ältere Anleitungen verwenden \cmd{git checkout -b ads}.
Das bewirkt dasselbe; \cmd{checkout} wurde aufgeteilt, weil es zu viele verschiedene
Aufgaben hatte.
\end{hintergrund}

% =====================================================================
\befehl{git branch}{Branches auflisten und löschen}

\begin{lstlisting}[style=befehle]
git branch                             # lokale Branches; * markiert den aktuellen
git branch -a                          # auch die Branches auf GitHub (all)
git branch -d ads                      # löschen (delete), nur wenn bereits übernommen
\end{lstlisting}

\cmd{-d} ist bewusst vorsichtig und verweigert das Löschen eines Branches, dessen Arbeit
noch nicht übernommen wurde. Das grosse \cmd{-D} erzwingt es und kann Arbeit kosten.

% =====================================================================
\befehl{git push}{Commits hochladen}

\textbf{Zweck:} schickt lokale Commits zum Remote.
\textbf{Weg:} lokales Repository $\rightarrow$ GitHub.

\begin{lstlisting}[style=befehle]
git push -u origin ads
\end{lstlisting}

\begin{teile}
  origin & das Ziel: der Name, unter dem GitHub als Remote eingetragen ist \\
  ads    & der Branch, der hochgeladen wird \\
  -u     & kurz für \cmd{--set-upstream}: verknüpft den lokalen Branch dauerhaft mit dem auf GitHub \\
\end{teile}

Nach dem ersten Push mit \cmd{-u} genügt \cmd{git push} ohne weitere Angaben. Hochgeladen
werden ausschliesslich \textbf{Commits}; Vorgemerktes und Ungespeichertes bleibt lokal.

\subsection*{Die Ausgabe lesen}

\begin{center}
\renewcommand{\arraystretch}{1.25}
\begin{tabularx}{\linewidth}{@{}>{\raggedright\arraybackslash\ttfamily\small}p{5.6cm}>{\raggedright\arraybackslash}X@{}}
  \toprule
  \normalfont\normalsize\textbf{Zeile} & \textbf{Bedeutung} \\
  \midrule
  Enumerating objects: 35 & Git ermittelt, welche gespeicherten Objekte zum Branch gehören \\
  Delta compression using up to 12 threads & überträgt nur Unterschiede statt ganzer Dateien, verteilt auf die Prozessorkerne \\
  Writing objects: 100\% (18/18), 2.29 KiB & tatsächlich übertragen: 18 Objekte, 2,29\,KB \\
  remote: \ldots & Zeilen mit \cmd{remote:} stammen von GitHub, nicht vom eigenen Computer \\
  {*} [new branch] ads -> ads & der Branch existiert jetzt auch auf GitHub \\
  branch \textquotesingle{}ads\textquotesingle{} set up to track \textquotesingle{}origin/ads\textquotesingle{} & Wirkung von \cmd{-u}: lokal und GitHub sind verknüpft \\
  \bottomrule
\end{tabularx}
\end{center}

% =====================================================================
\section{\texorpdfstring{\cmd{git fetch} und \cmd{git pull}}{git fetch und git pull}: Änderungen holen}

\textbf{Zweck:} holt Commits vom Remote.
\textbf{Weg:} GitHub $\rightarrow$ lokales Repository (und Arbeitsordner).

\begin{lstlisting}[style=befehle]
git fetch                              # nur herunterladen, nichts verändern
git pull                               # herunterladen und einbauen
\end{lstlisting}

\cmd{git pull} besteht aus zwei Schritten: Zuerst \cmd{git fetch}, das die neuen Commits
herunterlädt und \cmd{origin/main} aktualisiert, ohne die eigenen Dateien anzurühren. Danach
\cmd{git merge}, das die neuen Commits in den aktuellen Branch einbaut.

\begin{erfahrung}
\cmd{git fetch} ist immer ungefährlich. Wer unsicher ist, holt zuerst mit \cmd{fetch} und
schaut mit \cmd{git status}, ob der eigene Branch vor (\emph{ahead}) oder hinter
(\emph{behind}) GitHub liegt, bevor er mit \cmd{pull} etwas verändert.
\end{erfahrung}

% =====================================================================
\befehl{git stash}{Arbeit kurz beiseitelegen}

\textbf{Zweck:} räumt nicht committete Änderungen vorübergehend weg, damit der Arbeitsordner
sauber ist, zum Beispiel um kurz den Branch zu wechseln.

\begin{lstlisting}[style=befehle]
git stash                              # Änderungen beiseitelegen
git switch main                        # etwas anderes erledigen
git switch ads
git stash pop                          # Änderungen zurückholen
\end{lstlisting}

\begin{zusammenfassung}
\begin{itemize}[nosep]
  \item \cmd{status} und \cmd{diff} zeigen, verändern aber nichts.
  \item \cmd{add}, \cmd{commit}, \cmd{push} bewegen Änderungen nach rechts im Modell,
        \cmd{restore} und \cmd{pull} nach links.
  \item Optionen sind unabhängige Bausteine: einzeln ausprobieren, dann kombinieren.
\end{itemize}
\end{zusammenfassung}

% ---------------------------------------------------------------------
\chapter{Pull Requests und Code Review}
\label{kap:pr}

\begin{lernziele}
\begin{itemize}[nosep]
  \item verstehen, was ein Pull Request ist und warum er in der Industrie zentral ist,
  \item wissen, was bei einem Code Review geprüft wird,
  \item die drei Arten der Übernahme kennen und unterscheiden,
  \item Pull Requests mit \cmd{gh} erstellen und übernehmen.
\end{itemize}
\end{lernziele}

\section{Was ein Pull Request ist}

\begin{definition}{Pull Request}
Ein \begriff{Pull Request} (kurz PR) ist ein Antrag auf GitHub: „Bitte übernimm die Commits
meines Branches in \cmd{main}.“ Er bündelt die Änderungen, ihre Beschreibung, die
Diskussion darüber und die Ergebnisse automatischer Tests an einem Ort.
\end{definition}

Git selbst kennt keine Pull Requests. Sie sind eine Funktion der Plattform. Bei GitLab
heisst dasselbe \emph{Merge Request}.

Die Seite eines Pull Requests auf GitHub hat vier Reiter:

\begin{description}[style=nextline, leftmargin=1.2cm]
  \item[Conversation] Beschreibung, Kommentare und Freigaben.
  \item[Commits] die einzelnen Commits des Branches.
  \item[Checks] die Ergebnisse automatischer Tests.
  \item[Files changed] alle Änderungen als Diff. Hier findet das Review statt: Zu jeder
    Zeile kann man einen Kommentar schreiben.
\end{description}

\section{Code Review}

Beim \begriff{Code Review} liest mindestens eine zweite Person die Änderungen, bevor sie
übernommen werden. Typische Fragen dabei:

\begin{itemize}
  \item \textbf{Korrektheit:} Tut der Code, was er soll? Sind Sonderfälle bedacht, etwa
        leere Listen oder fehlende Eingaben?
  \item \textbf{Verständlichkeit:} Versteht jemand anderes den Code ohne Erklärung?
        Sind Namen sprechend?
  \item \textbf{Umfang:} Enthält der Pull Request nur, was zur Aufgabe gehört?
  \item \textbf{Sicherheit:} Stehen Passwörter im Code? Werden Eingaben geprüft?
  \item \textbf{Tests:} Ist die neue Funktion durch Tests abgesichert?
\end{itemize}

\begin{praxis}
Studien und Erfahrungen grosser Firmen zeigen, dass Reviews viele Fehler finden, bevor sie
Kundinnen und Kunden erreichen. Mindestens so wichtig ist aber ein Nebeneffekt: Wissen
verteilt sich im Team. Wenn nur eine Person einen Programmteil kennt und krank wird, steht
die Arbeit still. Durch Reviews kennen immer mindestens zwei Personen jede Änderung.

Darum gilt in Teams: \textbf{Kleine Pull Requests werden gründlich geprüft, grosse nur
überflogen.} Ein Pull Request sollte in 15 bis 30 Minuten verständlich sein.
\end{praxis}

\section{Automatische Prüfungen (CI)}

Bei jedem Pull Request kann GitHub automatisch Programme ausführen, zum Beispiel Tests,
Prüfungen des Codestils oder einen Testbuild. Das heisst \begriff{Continuous Integration}
(CI), auf GitHub umgesetzt mit \emph{GitHub Actions}. Die Anweisungen dafür liegen als
Datei im Repository unter \cmd{.github/workflows/}.

Das Ergebnis erscheint im Pull Request als grüner Haken oder rotes Kreuz. Rot bedeutet:
Etwas ist kaputt, die Übernahme wird blockiert. So gelangt fehlerhafter Code gar nicht erst
nach \cmd{main}.

\section{Drei Arten der Übernahme}

Beim Übernehmen eines Pull Requests bietet GitHub drei Varianten an. Ausgangslage: Der
Branch enthält drei Commits.

\begin{center}
\renewcommand{\arraystretch}{1.4}
\begin{tabularx}{\linewidth}{@{}>{\ttfamily}l>{\raggedright\arraybackslash}X>{\raggedright\arraybackslash}X@{}}
  \toprule
  \normalfont\textbf{Option} & \textbf{Ergebnis in main} & \textbf{Typischer Einsatz} \\
  \midrule
  --merge  & alle drei Commits bleiben erhalten, dazu ein \emph{Merge-Commit}, der die Übernahme festhält
           & wenn die einzelnen Schritte nachvollziehbar bleiben sollen; gut für den Einstieg \\
  --squash & die drei Commits werden zu \textbf{einem} zusammengefasst
           & sehr verbreitet in Firmen: \cmd{main} enthält genau einen Commit pro Aufgabe \\
  --rebase & die drei Commits werden einzeln an \cmd{main} angehängt, ohne Merge-Commit
           & wenn eine geradlinige Geschichte ohne Verzweigungen gewünscht ist \\
  \bottomrule
\end{tabularx}
\end{center}

\begin{erfahrung}
Es gibt keine „richtige“ Variante, nur eine einheitliche. Teams legen sich auf eine fest,
damit die Geschichte von \cmd{main} einheitlich aussieht. In diesem Handbuch wird
\cmd{--merge} verwendet, weil man im Log genau sieht, was wann übernommen wurde.
\end{erfahrung}

\section{Pull Requests auf der Kommandozeile}

\begin{lstlisting}[style=befehle]
gh pr create --fill
\end{lstlisting}

\begin{teile}
  pr create & Pull Request vom aktuellen Branch nach \cmd{main} erstellen \\
  --fill    & Titel und Beschreibung aus den Commit-Nachrichten übernehmen, ohne Rückfragen \\
\end{teile}

\begin{lstlisting}[style=befehle]
gh pr merge --merge --delete-branch
\end{lstlisting}

\begin{teile}
  pr merge        & den Pull Request des aktuellen Branches übernehmen \\
  --merge         & Art der Übernahme (siehe oben) \\
  --delete-branch & danach den Branch auf GitHub und lokal löschen und zu \cmd{main} wechseln \\
\end{teile}

Weitere nützliche Befehle:

\begin{lstlisting}[style=befehle]
gh pr list                             # offene Pull Requests anzeigen
gh pr view --web                       # aktuellen Pull Request im Browser öffnen
gh pr checks                           # Stand der automatischen Prüfungen
\end{lstlisting}

\section{Eine gute PR-Beschreibung}

Die Beschreibung richtet sich an die prüfende Person. Sie beantwortet drei Fragen:

\begin{lstlisting}[style=ausgabe]
## Was ändert sich?
Stack mit verketteter Liste statt Array implementiert.

## Warum?
Bei unbekannter Datenmenge braucht der Array-Stack teures Umkopieren.
Die Liste wächst in O(1) pro Element.

## Wie getestet?
Eigene Tests in stack_test.py, alle grün.
\end{lstlisting}

\begin{zusammenfassung}
\begin{itemize}[nosep]
  \item Ein Pull Request ist ein Antrag auf Übernahme mit Diskussion und Testergebnissen.
  \item Code Review findet Fehler und verteilt Wissen; kleine PRs werden besser geprüft.
  \item CI testet automatisch, rot blockiert die Übernahme.
  \item \cmd{--merge}, \cmd{--squash}, \cmd{--rebase}: Wichtig ist, dass ein Team sich auf eine Variante einigt.
\end{itemize}
\end{zusammenfassung}


\part{Praxis}
% ---------------------------------------------------------------------
\chapter{Gute Commits}
\label{kap:commits}

\begin{lernziele}
\begin{itemize}[nosep]
  \item wissen, wie gross ein Commit sein sollte und warum,
  \item Commit-Nachrichten nach gängiger Industrie-Konvention schreiben,
  \item die Konvention \emph{Conventional Commits} kennen.
\end{itemize}
\end{lernziele}

\section{Warum die Qualität von Commits zählt}

Ein Commit wird einmal geschrieben, aber oft gelesen: beim Review, bei der Fehlersuche
Monate später, von neuen Teammitgliedern. Die Geschichte eines Projekts ist seine
Dokumentation über das \emph{Warum}. Der Code zeigt, \emph{was} das Programm tut. Nur die
Commits erklären, \emph{weshalb} es so geworden ist.

\section{Ein Anliegen pro Commit}

\begin{definition}{Atomarer Commit}
Ein \begriff{atomarer Commit} enthält genau ein abgeschlossenes Anliegen, zum Beispiel eine
Funktion, eine Fehlerbehebung oder eine Umbenennung. Er lässt sich für sich allein
verstehen, prüfen und bei Bedarf rückgängig machen.
\end{definition}

\begin{center}
\renewcommand{\arraystretch}{1.3}
\begin{tabularx}{\linewidth}{@{}>{\raggedright\arraybackslash}X>{\raggedright\arraybackslash}X@{}}
  \toprule
  \textbf{Ungünstig} & \textbf{Besser} \\
  \midrule
  ein Commit „Abend 3“ mit Stack, Queue, Tippfehler in der README und neuer Ordnerstruktur
  & vier Commits: Stack, Queue, README korrigiert, Ordnerstruktur \\
  ein Commit pro gespeicherter Datei, auch halbfertig
  & ein Commit, sobald ein Teil funktioniert \\
  \bottomrule
\end{tabularx}
\end{center}

Stellt sich eine Woche später heraus, dass die Queue fehlerhaft ist, lässt sich im
zweiten Fall genau dieser eine Commit rückgängig machen, ohne Stack und README anzurühren.

\section{Die Commit-Nachricht}

Über Jahre hat sich ein Format etabliert, das von Werkzeugen wie GitHub und \cmd{git log}
optimal dargestellt wird:

\begin{lstlisting}[style=ausgabe]
ADS U02: Stack mit verketteter Liste implementiert

Der Array-Stack musste bei jeder Vergrösserung umkopieren. Die
verkettete Liste wächst pro Element in konstanter Zeit, was bei
unbekannter Datenmenge planbarer ist.
\end{lstlisting}

\begin{description}[style=nextline, leftmargin=1.2cm]
  \item[Betreffzeile] höchstens etwa 50 bis 70 Zeichen, beschreibt \emph{was} der Commit
    bewirkt. Sie erscheint in \cmd{git log --oneline} und als Titel auf GitHub.
  \item[Leerzeile] trennt Betreff und Text. Werkzeuge erkennen daran die Betreffzeile.
  \item[Text (optional)] erklärt \emph{warum}, wenn es nicht offensichtlich ist.
    Zeilen nach etwa 72 Zeichen umbrechen.
\end{description}

Mehrzeilige Nachrichten lassen sich auch ohne Editor schreiben: Jedes weitere \cmd{-m}
wird zu einem eigenen Absatz.

\begin{lstlisting}[style=befehle]
git commit -m "ADS U02: Stack mit verketteter Liste implementiert" -m "Der Array-Stack ..."
\end{lstlisting}

\subsection*{Gute und schlechte Beispiele}

\begin{center}
\renewcommand{\arraystretch}{1.3}
\begin{tabularx}{\linewidth}{@{}>{\ttfamily}l>{\raggedright\arraybackslash}X@{}}
  \toprule
  \normalfont\textbf{Nachricht} & \textbf{Beurteilung} \\
  \midrule
  update            & sagt nichts; was wurde aktualisiert, und warum? \\
  fix               & welcher Fehler? \\
  asdf              & leider häufig, völlig wertlos \\
  Änderungen        & beschreibt jeden denkbaren Commit \\
  \midrule
  LIAL SW03: LR-Zerlegung mit Pivotsuche     & klar, auffindbar, mit Bereich \\
  Spickzettel ADS: Heap-Operationen ergänzt  & klar, was und wo \\
  README: Einrichtung für Windows beschrieben & klar \\
  \bottomrule
\end{tabularx}
\end{center}

\section{Conventional Commits}

Viele Open-Source-Projekte und Firmen verwenden die Konvention \begriff{Conventional
Commits}. Jede Betreffzeile beginnt mit einem Typ:

\begin{center}
\renewcommand{\arraystretch}{1.25}
\begin{tabularx}{\linewidth}{@{}>{\ttfamily}l>{\raggedright\arraybackslash}X>{\ttfamily}l@{}}
  \toprule
  \normalfont\textbf{Typ} & \textbf{Bedeutung} & \normalfont\textbf{Beispiel} \\
  \midrule
  feat:     & neue Funktion          & feat: Export als CSV \\
  fix:      & Fehlerbehebung         & fix: Absturz bei leerer Liste \\
  docs:     & nur Dokumentation      & docs: Git-Handbuch Kapitel 9 \\
  refactor: & Umbau ohne neues Verhalten & refactor: Stack in eigene Klasse \\
  test:     & nur Tests              & test: Randfälle für Heap \\
  chore:    & Wartung, Konfiguration & chore: .gitignore ergänzt \\
  \bottomrule
\end{tabularx}
\end{center}

\begin{praxis}
Der Typ ist mehr als Ordnung: Werkzeuge lesen ihn automatisch aus. Aus allen
\cmd{feat:}- und \cmd{fix:}-Commits entstehen so die Versionshinweise einer neuen
Software-Version, ohne dass jemand sie von Hand schreibt.
\end{praxis}

\begin{erfahrung}
Ein guter Test für jede Commit-Nachricht: Ergibt der Satz „Wenn man diesen Commit
übernimmt, dann wird \ldots“ mit der Betreffzeile einen sinnvollen Satz? „\ldots{} dann wird
der Stack mit verketteter Liste implementiert“: ja. „\ldots{} dann wird update“: nein.
\end{erfahrung}

\begin{zusammenfassung}
\begin{itemize}[nosep]
  \item Ein Commit, ein Anliegen. Kleine Commits lassen sich prüfen und gezielt rückgängig machen.
  \item Betreffzeile kurz und konkret, Leerzeile, dann bei Bedarf das Warum.
  \item Conventional Commits (\cmd{feat:}, \cmd{fix:}, \ldots) ist ein verbreiteter Standard.
\end{itemize}
\end{zusammenfassung}

% ---------------------------------------------------------------------
\chapter{Zusammenarbeit und Konflikte}
\label{kap:zusammenarbeit}

\begin{lernziele}
\begin{itemize}[nosep]
  \item mit mehreren Geräten oder Personen am selben Repository arbeiten,
  \item verstehen, wie ein Konflikt entsteht,
  \item einen Konflikt Schritt für Schritt auflösen,
  \item Konflikte durch Arbeitsgewohnheiten vermeiden.
\end{itemize}
\end{lernziele}

\section{Mehrere Geräte, ein Repository}

PC und Laptop besitzen je einen eigenen Klon. Sie wissen nichts voneinander; verbunden sind
sie nur über GitHub:

\begin{center}
\begin{tikzpicture}[
    geraet/.style={draw=hsluBlau, fill=hsluBlau!6, rounded corners=3pt, thick,
                   minimum width=2.6cm, minimum height=1cm, font=\small\bfseries},
    weg/.style={-{Stealth[length=2.5mm]}, thick},
    lab/.style={font=\ttfamily\footnotesize, fill=white, inner sep=1pt}
  ]
  \node[geraet] (pc) at (0,0)    {PC};
  \node[geraet] (gh) at (5,0)    {GitHub};
  \node[geraet] (lt) at (10,0)   {Laptop};
  \draw[weg] ([yshift=2.5mm]pc.east)  -- node[lab, above=1mm] {push} ([yshift=2.5mm]gh.west);
  \draw[weg] ([yshift=-2.5mm]gh.west) -- node[lab, below=1mm] {pull} ([yshift=-2.5mm]pc.east);
  \draw[weg] ([yshift=2.5mm]gh.east)  -- node[lab, above=1mm] {pull} ([yshift=2.5mm]lt.west);
  \draw[weg] ([yshift=-2.5mm]lt.west) -- node[lab, below=1mm] {push} ([yshift=-2.5mm]gh.east);
\end{tikzpicture}
\end{center}

Mehrere Personen funktionieren genauso, nur dass auf jedem Gerät jemand anderes arbeitet.
Daraus folgt eine einfache Regel:

\begin{merke}
\textbf{Zu Beginn \cmd{git pull}, am Ende committen, pushen, mergen.} Wer das einhält, hat
auf allen Geräten denselben Stand und erlebt Konflikte nur noch selten.
\end{merke}

\subsection*{Was Git überträgt und was nicht}

\begin{itemize}
  \item \textbf{Mit Git übertragen} wird alles, was im Repository committet ist, auch
        Einstellungen wie \cmd{.vscode/settings.json} und Skripte wie \cmd{setup/einrichten.ps1}.
  \item \textbf{Nicht übertragen} werden installierte Programme, Anmeldungen wie die von
        \cmd{gh} und alles, was in \cmd{.gitignore} steht.
\end{itemize}

Darum enthält dieses Repository ein Einrichtungs-Skript. Auf jedem neuen Gerät genügt nach
dem Klonen ein Befehl:

\begin{lstlisting}[style=befehle]
powershell -ExecutionPolicy Bypass -File setup\einrichten.ps1
\end{lstlisting}

\begin{teile}
  powershell              & startet eine neue PowerShell, die das Skript ausführt \\
  -ExecutionPolicy Bypass & erlaubt Skripte nur für diesen Aufruf; Windows blockiert sie sonst standardmässig \\
  -File                   & die auszuführende Skript-Datei \\
\end{teile}

\begin{praxis}
Dieses Vorgehen ist in Firmen Standard und heisst oft \emph{Onboarding-Skript}: Wer neu
ins Team kommt, klont das Repository, startet ein Skript und kann am selben Tag
produktiv arbeiten. Was ein Gerät braucht, steht als Code im Repository statt in den
Köpfen einzelner Personen.
\end{praxis}

\section{Wie ein Konflikt entsteht}

Git führt Änderungen meist selbstständig zusammen, auch wenn zwei Personen dieselbe Datei
bearbeitet haben, solange es \emph{verschiedene Stellen} sind. Ein \begriff{Konflikt}
entsteht nur, wenn \textbf{dieselbe Stelle} auf zwei Wegen unterschiedlich geändert wurde.
Git kann dann nicht wissen, welche Fassung richtig ist, und fragt den Menschen.

Typisches Beispiel mit zwei Geräten:

\begin{enumerate}
  \item Am PC wird in \cmd{README.md} Zeile 5 geändert und gepusht.
  \item Am Laptop wurde vorher \emph{kein} \cmd{git pull} gemacht. Dort wird dieselbe Zeile
        anders geändert und committet.
  \item \cmd{git pull} am Laptop meldet: \cmd{CONFLICT (content): Merge conflict in README.md}.
\end{enumerate}

\section{Einen Konflikt auflösen}

Git markiert die betroffene Stelle direkt in der Datei:

\begin{lstlisting}[style=ausgabe]
<<<<<<< HEAD
Wöchentlich zwei Übungen lösen.
=======
Täglich eine Übung lösen.
>>>>>>> 7f02c1a (README: Lernrhythmus angepasst)
\end{lstlisting}

\begin{teile}
  <<<<<<< HEAD    & Beginn der eigenen Fassung, also des aktuellen Stands (hier: Laptop) \\
  =======         & Trennlinie zwischen den beiden Fassungen \\
  >>>>>>> 7f02c1a & Ende der anderen Fassung; dahinter Hash und Nachricht des Commits, der von GitHub kam \\
\end{teile}

Vorgehen in vier Schritten:

\begin{enumerate}
  \item Datei öffnen und die Stelle suchen. VS Code hebt sie farbig hervor und bietet
        Schaltflächen wie \emph{Accept Current} oder \emph{Accept Incoming}.
  \item Die gewünschte Endfassung herstellen: eine der beiden behalten oder beide
        sinnvoll kombinieren.
  \item \textbf{Alle Markierungszeilen} (\cmd{<<<<<<<}, \cmd{=======}, \cmd{>>>>>>>}) entfernen.
  \item Den aufgelösten Stand speichern und abschliessen:
\end{enumerate}

\begin{lstlisting}[style=befehle]
git add README.md                      # "Konflikt ist gelöst" melden
git commit -m "Konflikt in README aufgelöst"
git push
\end{lstlisting}

\begin{fehler}
Markierungszeilen werden übersehen und mitcommittet. Danach steht \cmd{<<<<<<< HEAD} im Code,
und das Programm startet nicht mehr. Vor dem \cmd{git add} darum die Datei nach
\cmd{<<<<<<<} durchsuchen (Strg+F).
\end{fehler}

\begin{erfahrung}
Ein Konflikt ist kein Fehler und nichts Gefährliches. Git hat nichts verloren, es fragt
lediglich nach. Wer unsicher ist, kann den Vorgang jederzeit mit \cmd{git merge --abort}
abbrechen; danach ist alles wie vor dem \cmd{pull}.
\end{erfahrung}

\section{Konflikte vermeiden}

\begin{itemize}
  \item Vor Arbeitsbeginn \cmd{git pull}.
  \item Häufig pushen, statt tagelang lokal zu sammeln.
  \item Kurzlebige Branches: je kürzer ein Branch lebt, desto weniger weicht er von \cmd{main} ab.
  \item Im Team absprechen, wer an welchen Dateien arbeitet.
\end{itemize}

\begin{zusammenfassung}
\begin{itemize}[nosep]
  \item Geräte und Personen sind nur über GitHub verbunden: am Anfang \cmd{pull}, am Ende \cmd{push}.
  \item Ein Konflikt entsteht nur, wenn dieselbe Stelle unterschiedlich geändert wurde.
  \item Auflösen: Stelle bereinigen, Markierungen löschen, \cmd{add}, \cmd{commit}, \cmd{push}.
  \item \cmd{git merge --abort} bricht ab und stellt den vorherigen Zustand wieder her.
\end{itemize}
\end{zusammenfassung}

% ---------------------------------------------------------------------
\chapter{Fehler beheben}
\label{kap:fehler}

\begin{lernziele}
\begin{itemize}[nosep]
  \item für jede typische Panne den passenden Befehl finden,
  \item verstehen, warum „schon gepusht?“ die entscheidende Frage ist,
  \item wissen, welche Befehle gefährlich sind und welche ein Sicherheitsnetz bieten.
\end{itemize}
\end{lernziele}

\section{Die entscheidende Frage: schon gepusht?}

Solange ein Commit nur lokal existiert, gehört er niemandem sonst. Man darf ihn ändern oder
ersetzen. Sobald er gepusht ist, haben andere ihn womöglich schon geholt und darauf
aufgebaut. Ein nachträgliches Ändern würde deren Geschichte widersprechen (siehe Hintergrund
zum Hash in Kapitel~\ref{kap:modell}). Daraus folgt die Grundregel:

\begin{merke}
\textbf{Lokal:} Commits dürfen korrigiert und ersetzt werden.\\
\textbf{Gepusht:} Commits werden nicht mehr verändert, sondern durch \emph{neue} Commits korrigiert.
\end{merke}

\section{Pannen und ihre Lösung}

\begin{center}
\renewcommand{\arraystretch}{1.4}
\small
\begin{tabularx}{\linewidth}{@{}>{\raggedright\arraybackslash}X>{\raggedright\arraybackslash\ttfamily}p{5.3cm}>{\raggedright\arraybackslash}X@{}}
  \toprule
  \textbf{Situation} & \normalfont\textbf{Befehl} & \textbf{Wirkung} \\
  \midrule
  falsche Datei vorgemerkt & git restore --staged <datei> & aus der Staging Area nehmen, Datei unverändert \\
  Änderungen an einer Datei verwerfen & git restore <datei> & zurück auf den letzten Commit, \textbf{endgültig} \\
  versehentlich auf \cmd{main} gearbeitet, noch nicht committet & git switch -c <name> & Änderungen wandern auf den neuen Branch \\
  Commit-Nachricht falsch, \textbf{nicht gepusht} & git commit --amend -m "neu" & ersetzt den letzten Commit \\
  Datei im Commit vergessen, \textbf{nicht gepusht} & git add <datei> \newline git commit --amend --no-edit & fügt sie dem letzten Commit hinzu; \cmd{--no-edit} behält die Nachricht \\
  fehlerhafter Commit, \textbf{bereits gepusht} & git revert <hash> & neuer Commit, der die Änderungen des alten aufhebt \\
  Arbeit kurz beiseitelegen & git stash \newline git stash pop & parkt und holt Änderungen zurück \\
  \bottomrule
\end{tabularx}
\end{center}

\subsection*{\texorpdfstring{\cmd{--amend}}{--amend} und \texorpdfstring{\cmd{revert}}{revert} im Vergleich}

\begin{description}[style=nextline, leftmargin=1.2cm]
  \item[\cmd{git commit --amend}] erzeugt einen neuen Commit, der den letzten
    \emph{ersetzt}. Der alte verschwindet aus der Geschichte. Das ist nur sicher, solange
    niemand den alten Commit kennt, also vor dem Push.
  \item[\cmd{git revert <hash>}] erzeugt einen \emph{zusätzlichen} Commit, der genau das
    Gegenteil des angegebenen Commits enthält. Die Geschichte bleibt vollständig und
    ehrlich: Man sieht den Fehler und seine Korrektur. Darum ist das der Weg für
    bereits geteilte Commits.
\end{description}

\begin{praxis}
Wenn in einer Firma eine fehlerhafte Änderung live gegangen ist, wird sie mit einem
\cmd{revert} zurückgenommen. Das geht schnell, ist nachvollziehbar und durchläuft denselben
Prüfprozess wie jede andere Änderung. Erst danach sucht man in Ruhe nach der eigentlichen
Ursache.
\end{praxis}

\section{Das Sicherheitsnetz: \texorpdfstring{\cmd{git reflog}}{git reflog}}

Git protokolliert jede Bewegung von HEAD, also jeden Commit, jeden Branch-Wechsel und jedes
Zurücksetzen. Dieses Protokoll heisst \begriff{Reflog} und lässt sich mit
\cmd{git reflog} anzeigen:

\begin{lstlisting}[style=ausgabe]
3b1a03a HEAD@{0}: pull: Fast-forward
db753be HEAD@{1}: commit: ADS: Uebungen U01-U06
19ce440 HEAD@{2}: clone: from https://github.com/...
\end{lstlisting}

Selbst ein scheinbar verlorener Commit ist darin meist noch zu finden. Mit seinem Hash lässt
er sich zurückholen, zum Beispiel mit \cmd{git switch -c rettung <hash>}.

\begin{erfahrung}
In Git geht erstaunlich wenig wirklich verloren, solange es einmal committet war. Nicht
committete Änderungen sind dagegen ungeschützt. Darum lieber einmal zu oft committen,
zum Beispiel als Zwischenstand auf dem eigenen Branch.
\end{erfahrung}

\section{Gefährliche Befehle}

\begin{achtung}
Zwei Befehle können tatsächlich Arbeit vernichten:
\begin{description}[style=nextline, leftmargin=1.2cm]
  \item[\cmd{git reset --hard}] setzt Arbeitsordner und Staging Area hart auf einen
    früheren Stand zurück. Nicht committete Änderungen sind danach unwiederbringlich verloren.
  \item[\cmd{git push --force}] überschreibt die Geschichte auf GitHub mit der eigenen. Was
    andere inzwischen gepusht haben, verschwindet dort.
\end{description}
Beide nur verwenden, wenn man genau weiss, was sie bewirken. In Teams ist \cmd{--force}
auf \cmd{main} durch Branch-Schutz meist technisch gesperrt.
\end{achtung}

\begin{zusammenfassung}
\begin{itemize}[nosep]
  \item Vor dem Push darf korrigiert werden (\cmd{--amend}), danach wird mit neuen Commits korrigiert (\cmd{revert}).
  \item \cmd{git restore --staged} ist harmlos, \cmd{git restore} ohne Option verwirft endgültig.
  \item \cmd{git reflog} findet fast jeden committeten Stand wieder.
  \item \cmd{reset --hard} und \cmd{push --force} nur mit vollem Verständnis.
\end{itemize}
\end{zusammenfassung}

% ---------------------------------------------------------------------
\chapter{Meldungen verstehen}
\label{kap:meldungen}

\begin{lernziele}
\begin{itemize}[nosep]
  \item die häufigsten Meldungen von Git und \cmd{gh} einordnen,
  \item zwischen harmlosen Hinweisen und echten Fehlern unterscheiden.
\end{itemize}
\end{lernziele}

Git unterscheidet drei Arten von Meldungen. Sie beginnen meist mit einem Schlüsselwort:

\begin{center}
\renewcommand{\arraystretch}{1.25}
\begin{tabularx}{\linewidth}{@{}>{\ttfamily}l>{\raggedright\arraybackslash}X@{}}
  \toprule
  \normalfont\textbf{Beginnt mit} & \textbf{Bedeutung} \\
  \midrule
  hint:    & Tipp; der Befehl war erfolgreich oder Git schlägt den nächsten Schritt vor \\
  warning: & Hinweis; der Befehl wurde ausgeführt, aber etwas ist bemerkenswert \\
  error: / fatal: & Fehler; der Befehl wurde \textbf{nicht} ausgeführt \\
  \bottomrule
\end{tabularx}
\end{center}

\section{Zustandsmeldungen}

\meldung{nothing to commit, working tree clean}
  {Es gibt keine Änderungen. Arbeitsordner und letzter Commit sind identisch.}
  {Nichts. Alles ist gespeichert.}

\meldung{Your branch is ahead of \textquotesingle{}origin/main\textquotesingle{} by 1 commit.}
  {Ein lokaler Commit ist noch nicht auf GitHub.}
  {\cmd{git push}}

\meldung{Your branch is behind \textquotesingle{}origin/main\textquotesingle{} by 2 commits.}
  {Auf GitHub gibt es neue Commits, die lokal fehlen.}
  {\cmd{git pull}}

\section{Hinweise}

\meldung{warning: in the working copy of \textquotesingle{}datei.py\textquotesingle, LF will be replaced by CRLF}
  {Die Datei verwendet Zeilenenden im Linux-Format. Git wird sie unter Windows anpassen
   (siehe Abschnitt~\ref{sec:zeilenenden}).}
  {Nichts. Reine Information.}

\section{Fehler}

\meldung{error: src refspec ads does not match any}
  {\cmd{git push} soll den Branch \cmd{ads} hochladen, der lokal nicht existiert.
   Meist wurde \cmd{git switch -c ads} vergessen und auf \cmd{main} committet.}
  {Mit \cmd{git branch} prüfen. Den Branch nachträglich anlegen: \cmd{git switch -c ads},
   dann erneut pushen.}

\meldung{no pull requests found for branch "ads"}
  {\cmd{gh pr merge} findet keinen Pull Request, weil \cmd{gh pr create} übersprungen wurde.}
  {Zuerst \cmd{gh pr create --fill}, dann mergen.}

\meldung{fatal: not a git repository (or any of the parent directories): .git}
  {Das Terminal steht in einem Ordner, der zu keinem Repository gehört.}
  {Mit \cmd{cd} in den Repository-Ordner wechseln.}

\meldung{! [rejected] main -> main (fetch first)}
  {Auf GitHub gibt es Commits, die lokal fehlen. Git verweigert das Überschreiben, um
   fremde Arbeit zu schützen.}
  {Zuerst \cmd{git pull}, gegebenenfalls Konflikte lösen, dann \cmd{git push}.}

\meldung{CONFLICT (content): Merge conflict in README.md}
  {Dieselbe Stelle wurde auf zwei Wegen unterschiedlich geändert.}
  {Siehe Kapitel~\ref{kap:zusammenarbeit}.}

\meldung{could not prompt: Incorrect function. You appear to be running in MinTTY}
  {Git Bash kann die Auswahlmenüs von \cmd{gh} nicht darstellen.}
  {Alle Antworten als Optionen mitgeben oder PowerShell verwenden.}

\meldung{refusing to allow an OAuth App to create or update workflow \textasciigrave.github/workflows/ci.yml\textasciigrave{} without \textasciigrave workflow\textasciigrave{} scope}
  {Der Push enthält eine Datei unter \cmd{.github/workflows/}. Die Anmeldung von \cmd{gh}
   hat dafür keine Berechtigung, weil Workflows automatisch Code ausführen.}
  {\cmd{gh auth refresh -h github.com -s workflow} und im Browser bestätigen.}

\section{Zeilenenden}
\label{sec:zeilenenden}

Ein Zeilenumbruch ist in einer Textdatei ein unsichtbares Zeichen. Linux und macOS
verwenden eines, \emph{Line Feed} (LF). Windows verwendet zwei, \emph{Carriage Return} und
\emph{Line Feed} (CRLF). Die Namen stammen von der Schreibmaschine: Wagen zurückschieben,
Papier eine Zeile weiterdrehen.

Mit der Einstellung \cmd{core.autocrlf = true} wandelt Git beim Auschecken unter Windows in
CRLF um und speichert im Repository LF. So passt jede Datei auf jedem System, und
Unterschiede in den Zeilenenden erscheinen nicht als Änderungen.

\begin{zusammenfassung}
\begin{itemize}[nosep]
  \item \cmd{hint:} und \cmd{warning:} sind Informationen, \cmd{error:} und \cmd{fatal:} echte Fehler.
  \item Fast jede Fehlermeldung nennt Ursache und Lösung. In Ruhe lesen lohnt sich.
\end{itemize}
\end{zusammenfassung}

% ---------------------------------------------------------------------
\chapter{Sicherheit und Ordnung im Repository}
\label{kap:sicherheit}

\begin{lernziele}
\begin{itemize}[nosep]
  \item wissen, was nicht in ein Repository gehört und warum,
  \item eine \cmd{.gitignore} lesen und schreiben,
  \item richtig reagieren, wenn ein Passwort versehentlich auf GitHub gelandet ist.
\end{itemize}
\end{lernziele}

\section{Was nicht ins Repository gehört}

\begin{center}
\renewcommand{\arraystretch}{1.35}
\begin{tabularx}{\linewidth}{@{}>{\raggedright\arraybackslash}p{4.2cm}>{\raggedright\arraybackslash}X>{\raggedright\arraybackslash}p{3.6cm}@{}}
  \toprule
  \textbf{Was} & \textbf{Warum nicht} & \textbf{Stattdessen} \\
  \midrule
  Passwörter, API-Schlüssel, Tokens & werden mit dem Repository geteilt und bleiben in der Geschichte & Datei \cmd{.env}, in \cmd{.gitignore} \\
  erzeugte Dateien (PDF aus LaTeX, \cmd{__pycache__}, \cmd{build/}) & lassen sich jederzeit neu erzeugen und verursachen unnötige Konflikte & in \cmd{.gitignore} \\
  grosse Dateien (Videos, Datensätze) & Git speichert jede Version vollständig; das Repository wird riesig. GitHub lehnt Dateien über 100\,MB ab & Cloud-Speicher, Git LFS \\
  persönliche Einstellungen, Daten von Programmen & betreffen nur ein Gerät & lokal lassen \\
  urheberrechtlich geschütztes Material Dritter & in öffentlichen Repositorys eine Rechtsverletzung & privates Repository oder Verweis \\
  \bottomrule
\end{tabularx}
\end{center}

\begin{praxis}
Die Faustregel in der Industrie lautet: \textbf{Ins Repository gehört die Quelle, nicht das
Ergebnis.} Die \cmd{.tex}-Datei kommt hinein, das daraus erzeugte PDF nicht. Der
Python-Code kommt hinein, die erzeugten \cmd{.pyc}-Dateien nicht. Jeder kann das Ergebnis
aus der Quelle neu herstellen.
\end{praxis}

\section{Die Datei .gitignore}

Die Datei \cmd{.gitignore} im Hauptordner listet Dateien und Ordner, die Git ignorieren soll.
Sie erscheinen dann nicht in \cmd{git status} und werden von \cmd{git add -A} übergangen.

\begin{lstlisting}[style=ausgabe]
# Python
__pycache__/
*.pyc
.venv/

# LaTeX: PDF entsteht beim Kompilieren neu
build/
*.aux
*.log

# Geheimnisse
.env
\end{lstlisting}

\begin{center}
\renewcommand{\arraystretch}{1.25}
\begin{tabularx}{\linewidth}{@{}>{\ttfamily}l>{\raggedright\arraybackslash}X@{}}
  \toprule
  \normalfont\textbf{Muster} & \textbf{Wirkung} \\
  \midrule
  build/       & der Ordner \cmd{build} an jeder Stelle, mit allem Inhalt \\
  *.log        & alle Dateien mit der Endung \cmd{.log}; \cmd{*} steht für beliebige Zeichen \\
  .env         & genau diese Datei \\
  !wichtig.log & Ausnahme: diese Datei trotz \cmd{*.log} nicht ignorieren \\
  \texttt{\#} Text & Kommentar \\
  \bottomrule
\end{tabularx}
\end{center}

Prüfen, warum eine Datei ignoriert wird:

\begin{lstlisting}[style=befehle]
git check-ignore -v build/ads.pdf      # -v (verbose): nennt Datei und Zeile der Regel
\end{lstlisting}

\begin{fehler}
Eine Datei wird nachträglich in \cmd{.gitignore} eingetragen, erscheint aber weiter in
\cmd{git status}. Ursache: \cmd{.gitignore} wirkt nur auf Dateien, die Git noch nicht kennt.
Eine bereits committete Datei muss zuerst aus dem Index genommen werden:
\cmd{git rm --cached <datei>}. Die Datei bleibt auf der Festplatte erhalten.
\end{fehler}

\section{Wenn ein Geheimnis auf GitHub gelandet ist}

\begin{achtung}
Ein Passwort oder Schlüssel, der einmal auf GitHub war, gilt als \textbf{veröffentlicht}.
Automatische Programme durchsuchen öffentliche Repositorys rund um die Uhr nach solchen
Schlüsseln, oft innerhalb von Minuten nach dem Push. Das Löschen der Datei in einem neuen
Commit genügt nicht: Der Schlüssel bleibt in der Geschichte lesbar.
\end{achtung}

Richtiges Vorgehen, in dieser Reihenfolge:

\begin{enumerate}
  \item \textbf{Den Schlüssel sofort ungültig machen} und einen neuen erzeugen, beim
        jeweiligen Dienst. Das ist der einzige Schritt, der den Schaden wirklich begrenzt.
  \item Die Datei aus dem Repository entfernen und in \cmd{.gitignore} eintragen.
  \item Bei Bedarf die Geschichte bereinigen lassen. Das ist aufwendig und gehört in
        erfahrene Hände.
\end{enumerate}

\section{Öffentlich oder privat?}

Ein \textbf{öffentliches} Repository kann jeder Mensch lesen und kopieren. Es eignet sich als
Portfolio für eigene Arbeiten und wird von Arbeitgebern gerne angesehen. Ein
\textbf{privates} Repository sehen nur eingeladene Personen. Dorthin gehören Unterlagen
Dritter, persönliche Notizen und alles, was nicht für die Öffentlichkeit bestimmt ist.

\begin{erfahrung}
Ein öffentliches GitHub-Profil ist heute oft Teil einer Bewerbung. Saubere Commits,
verständliche README-Dateien und eine nachvollziehbare Geschichte sagen über die Arbeitsweise
mehr aus als jedes Zeugnis.
\end{erfahrung}

\begin{zusammenfassung}
\begin{itemize}[nosep]
  \item Quelle ja, Ergebnis nein. Keine Geheimnisse, keine grossen Dateien.
  \item \cmd{.gitignore} steuert, was Git übersieht; bereits bekannte Dateien mit \cmd{git rm --cached} lösen.
  \item Geleakter Schlüssel: zuerst ungültig machen, dann aufräumen.
\end{itemize}
\end{zusammenfassung}


\appendix
% ---------------------------------------------------------------------
\chapter{Kurzreferenz}
\label{kap:kurzreferenz}

\section*{Täglicher Ablauf}

\begin{lstlisting}[style=befehle]
git switch main                       # 1. aufs Original wechseln
git pull                              #    und aktuellen Stand holen
git switch -c <name>                  # 2. Branch anlegen
git status                            # 3. Lage prüfen
git add <pfad>                        # 4. vormerken
git diff --cached --stat              # 5. Auswahl kontrollieren
git commit -m "<Bereich>: <was>"      # 6. Commit
git push -u origin <name>             # 7. hochladen
gh pr create --fill                   # 8. Pull Request
gh pr merge --merge --delete-branch   # 9. übernehmen
git pull                              # 10. Stand holen
\end{lstlisting}

\section*{Nach Aufgabe}

\begin{center}
\renewcommand{\arraystretch}{1.3}
\begin{tabularx}{\linewidth}{@{}>{\raggedright\arraybackslash}X>{\raggedright\arraybackslash\ttfamily}p{8.0cm}@{}}
  \toprule
  \textbf{Ich will \ldots} & \normalfont\textbf{Befehl} \\
  \midrule
  \multicolumn{2}{@{}l}{\textit{Lage prüfen}} \\
  sehen, wo ich stehe                 & git status \quad\textrm{oder}\quad git status -sb \\
  Vormerkung prüfen                   & git diff --cached --stat \\
  Änderungen im Detail sehen          & git diff \quad\textrm{bzw.}\quad git diff --cached \\
  Geschichte ansehen                  & git log --oneline --graph --all -10 \\
  Branches auflisten                  & git branch -a \\
  \midrule
  \multicolumn{2}{@{}l}{\textit{Speichern und teilen}} \\
  vormerken                           & git add <pfad> \\
  Commit machen                       & git commit -m "Nachricht" \\
  Branch anlegen                      & git switch -c <name> \\
  hochladen, erstes Mal               & git push -u origin <name> \\
  hochladen, danach                   & git push \\
  Pull Request erstellen              & gh pr create --fill \\
  Pull Request übernehmen             & gh pr merge --merge --delete-branch \\
  \midrule
  \multicolumn{2}{@{}l}{\textit{Holen}} \\
  neuesten Stand holen                & git pull \\
  nur nachsehen, ob es Neues gibt     & git fetch \newline git status \\
  \midrule
  \multicolumn{2}{@{}l}{\textit{Korrigieren}} \\
  aus Vormerkung nehmen               & git restore --staged <datei> \\
  Änderungen verwerfen (endgültig)    & git restore <datei> \\
  letzten Commit ändern (nicht gepusht) & git commit --amend \\
  gepushten Commit aufheben           & git revert <hash> \\
  Arbeit parken                       & git stash \newline git stash pop \\
  verlorenen Commit finden            & git reflog \\
  \midrule
  \multicolumn{2}{@{}l}{\textit{Einrichten}} \\
  Identität setzen                    & git config --global user.name "Name" \newline git config --global user.email "mail" \\
  bei GitHub anmelden                 & gh auth login --hostname github.com \newline \hspace*{1em}--git-protocol https --web \\
  Repository klonen                   & gh repo clone <besitzer>/<repo> \\
  Hilfe                               & git <befehl> -h \quad\textrm{oder}\quad git help <befehl> \\
  \bottomrule
\end{tabularx}
\end{center}

% ---------------------------------------------------------------------
\chapter{Glossar}
\label{kap:glossar}

\begin{description}[style=nextline, leftmargin=1.2cm, itemsep=3pt]
  \item[Branch \normalfont(Zweig)] eigene Arbeitslinie; technisch ein beweglicher Zeiger auf
    einen Commit. Siehe Kapitel~\ref{kap:modell}.
  \item[CI \normalfont(Continuous Integration)] automatische Tests und Prüfungen bei jeder
    Änderung, auf GitHub mit \emph{GitHub Actions}.
  \item[Clone \normalfont(Klon)] vollständige Kopie eines Repositorys inklusive Geschichte.
  \item[Code Review] Prüfung einer Änderung durch eine zweite Person vor der Übernahme.
  \item[Commit] gespeicherter Zwischenstand des ganzen Projekts mit Autor, Zeit, Nachricht
    und Verweis auf den Vorgänger.
  \item[Conventional Commits] verbreitete Konvention für Commit-Nachrichten mit Typen wie
    \cmd{feat:} und \cmd{fix:}.
  \item[Diff] Darstellung der Unterschiede zwischen zwei Ständen, Zeile für Zeile.
  \item[Fetch] Commits vom Remote herunterladen, ohne die eigenen Dateien zu verändern.
  \item[Git] verteiltes Versionsverwaltungssystem, 2005 von Linus Torvalds entwickelt.
  \item[GitHub] Internetdienst für Git-Repositorys mit Pull Requests, Code Review und CI.
  \item[Hash] eindeutige Kennung eines Commits aus 40 Zeichen; meist werden die ersten sieben
    verwendet.
  \item[HEAD] Zeiger auf den aktuell ausgecheckten Branch beziehungsweise Commit.
  \item[Index] anderer Name für die Staging Area.
  \item[Konflikt \normalfont(merge conflict)] dieselbe Stelle wurde auf zwei Wegen
    unterschiedlich geändert; ein Mensch muss entscheiden.
  \item[main] Hauptbranch mit dem gültigen Stand; früher oft \cmd{master}.
  \item[Merge] Zusammenführen zweier Arbeitslinien.
  \item[Merge-Commit] Commit mit zwei Vorgängern, der eine Zusammenführung festhält.
  \item[origin] üblicher Name des Remotes, von dem geklont wurde.
  \item[Pull] \emph{fetch} und \emph{merge} in einem Schritt.
  \item[Pull Request \normalfont(PR)] Antrag auf GitHub, die Commits eines Branches zu
    übernehmen; bei GitLab \emph{Merge Request}.
  \item[Push] eigene Commits zum Remote hochladen.
  \item[Rebase] Commits auf einen anderen Ausgangspunkt verschieben; ergibt eine geradlinige
    Geschichte.
  \item[Reflog] Protokoll aller Bewegungen von HEAD; Sicherheitsnetz zum Wiederfinden von
    Commits.
  \item[Remote] entfernte Kopie eines Repositorys, meist auf GitHub.
  \item[Repository \normalfont(Repo)] Projekt samt vollständiger Geschichte; liegt im Ordner
    \cmd{.git}.
  \item[Revert] neuer Commit, der die Änderungen eines früheren Commits aufhebt.
  \item[Squash] mehrere Commits zu einem zusammenfassen.
  \item[Staging Area \normalfont(Vormerkung)] Zwischenablage für den nächsten Commit.
  \item[Stash] vorübergehendes Versteck für nicht committete Änderungen.
  \item[Token] digitaler Schlüssel, mit dem ein Programm im Namen eines Benutzers handeln darf.
  \item[Untracked] Datei, die Git noch nie gesehen hat.
  \item[Upstream] der mit einem lokalen Branch verknüpfte Branch auf dem Remote.
  \item[Working Tree \normalfont(Arbeitsordner)] die sichtbaren Dateien, in denen gearbeitet wird.
\end{description}

\end{document}
